\documentclass[10pt,a4paper,draft]{article}

\usepackage{fullpage}
\usepackage{xparse} % provides \NewDocumentEnvironment
\usepackage{amsthm}
\usepackage{amsmath}
\usepackage{amssymb} % provides \lrcorner
\usepackage[all,cmtip]{xy} % provides \xymatrix
\usepackage[T1]{fontenc}    % use accented characters from the font
%\usepackage[utf8]{inputenc} % preloaded (>= 2018), provides \DeclareUnicodeCharacter
%\usepackage[notext]{stix2}  % provides \similarrightarrow, not contained in arXiv's TeXLive
\usepackage{textgreek} % "\textlambda-calculus"
\usepackage{enumitem} % provides \newlist
\usepackage[inline,nomargin]{fixme} % fixme notes
\fxsetup{targetlayout=color}
%\usepackage{appendix} % provides \appendixpage
%\usepackage{apptools} % provides \AtAppendix
\usepackage{chngcntr} % provides \counterwithout
\usepackage{hyphenat} % provides \hyp ("typesets a hyphen and allows full automatic hyphenation of the other words forming [a] compound word")
\usepackage{xspace}   % provides \xspace
\usepackage{mathtools} % provides \DeclarePairedDelimiter
\usepackage[final]{hyperref} % option final needed to get links with document operation "draft"
\usepackage[capitalise,nameinlink]{cleveref} % "cleveref must be loaded last (after other packages that modify LaTeX's referencing system)"

%%% Macros
\def\Type{\hbox{\sf Type}}
\newcommand{\nats}{\mathbb{N}}
\newcommand{\ints}{\mathbb{Z}}
\newcommand{\II}{\mathbb{I}}
\newcommand{\JJ}{\mathsf{J}}
\newcommand{\Empty}{\mathbf{0}}
\newcommand{\Unit}{\mathbf{1}}
\newcommand{\Nat}{\mathsf{Nat}}
\newcommand{\W}{\mathsf{W}}
\newcommand{\Id}{\mathsf{Id}}
\newcommand{\north}{\mathsf{north}}
\newcommand{\susp}{\mathop{\mathsf{Susp}}}
\newcommand{\south}{\mathsf{south}}
\newcommand{\merid}{\mathsf{merid}}
\newcommand{\Path}{{\sf Path}}
\def\norm#1{\left\|#1\right\|}
\newcommand{\UU}{\mathcal{U}}%{\mathscr{U}}%{\mathsf{U}}
\newcommand{\HProp}{\mathsf{HProp}}
\newcommand{\El}{\mathsf{El}}
\newcommand{\In}{\mathsf{In}}
\newcommand{\VV}{\mathcal{S}}%{\mathscr{S}}%{\mathsf{S}}
\newcommand{\Elem}{{\sf Elem}}
\newcommand{\Cov}{\mathsf{C}}%{\mathsf{Cov}}
\newcommand{\CC}{{\mathcal C}}
\newcommand{\GG}{\mathsf{isMod}}
\newcommand{\munit}{\inc}
\newcommand{\inc}{\mathsf{inc}}
\newcommand{\patch}{\mathsf{patch}}
\newcommand{\leftinv}{\mathsf{linv}}
\newcommand{\ext}{\mathsf{ext}}
\newcommand{\msup}{\mathsf{sup}}
\newcommand{\refl}{\mathsf{refl}}
\newcommand{\J}{\mathsf{J}}
\newcommand{\hrefl}{\mathsf{hrefl}}
\newcommand{\id}{\mathsf{id}}
\newcommand{\diag}{\mathsf{diag}}
\newcommand{\const}{\mathsf{const}}
\newcommand{\gap}{\mathsf{gap}}
\newcommand{\ap}{\mathsf{ap}}
\newcommand{\total}{\mathsf{total}}
\newcommand{\hFiber}{\mathsf{fib}}
\newcommand{\isContr}{\mathsf{isContr}}
\newcommand{\isEquiv}{\mathsf{isEquiv}}
\newcommand{\BB}{\mathcal{B}}
\newcommand{\MP}{\mathsf{MP}}
\newcommand{\FT}{\mathsf{FT}}
\newcommand{\pp}{{\sf p}}
\newcommand{\qq}{{\sf q}}
\newcommand{\unit}{()} % \mathbf{0}}
\newcommand{\true}{\mathsf{true}}
\newcommand{\false}{\mathsf{false}}
\newcommand{\zero}{\mathsf{zero}}
\newcommand{\SUCC}{\mathsf{succ}}
\newcommand{\natrec}{\mathsf{natrec}}
\newcommand{\intro}{\mathsf{intro}}
\newcommand{\rec}{\mathsf{rec}}
\newcommand{\elim}{\mathsf{elim}}
\newcommand{\wsup}{\mathsf{sup}}
\newcommand{\truncinc}{\inc}
\newcommand{\squash}{\mathsf{squash}}
\newcommand{\truncext}{\ext}
\newcommand{\lam}[1]{λ{#1}}
\newcommand{\abs}[1]{\lam{#1}}
\newcommand{\app}[2]{{#1}\;{#2}}
\newcommand{\fst}{\mathsf{fst}}
\newcommand{\snd}{\mathsf{snd}}
\newcommand{\pair}[2]{(#1,#2)}
\newcommand{\reflF}{\mathsf{refl}}
%\newcommand{\unitF}{\mathbf{\ast}}
%\newcommand{\zeroF}{\langle \rangle}
\newcommand{\unitF}{\langle \rangle}
\newcommand{\fstF}{\pi_1}
\newcommand{\sndF}{\pi_2}
\newcommand{\pairF}[2]{\langle #1,#2 \rangle}
\newcommand{\mpF}[1]{\langle #1 \rangle} % n-ary composition, \ncirc?
\newcommand{\D}{D}              % action on contexts
%\newcommand{\DD}{\widetilde{D}} % action on types
%\newcommand{\dd}{\widetilde{D}} % action on terms
\newcommand{\DD}{\tilde{D}} % action on types
\newcommand{\dd}{\tilde{D}} % action on terms
\newcommand{\point}{\eta}
\newcommand{\ppoint}{\widetilde{\eta}}
\newcommand{\E}{E}
%\newcommand{\EE}{\widetilde{D}}
%\newcommand{\ee}{\widetilde{d}}
\newcommand{\EE}{\tilde{E}}
\newcommand{\ee}{\tilde{e}}
\newcommand{\FF}{\widetilde{F}}
\renewcommand{\~}{\widetilde}
\newcommand{\PP}{\mathsf{P}}
\newcommand{\QQ}{\mathsf{Q}}
\newcommand{\RR}{\mathsf{R}}
\newcommand{\ua}{\mathsf{ua}}
% categories with families
\newcommand{\Ctx}{\mathsf{Ctx}}
\newcommand{\Ty}{Ty}
\newcommand{\Tm}{Tm}
\newcommand{\p}{\mathsf{p}}
\newcommand{\q}{\mathsf{q}}
\newcommand{\ir}[1]{[#1]}
\newcommand{\il}[1]{⟨#1⟩}
\newcommand{\term}{\mathsf{!}}
\newcommand{\Agda}{\textsc{Agda}\xspace}
\newcommand{\CIC}{\textsf{CIC}\xspace}

\newcommand{\op}{\mathsf{op}}
\DeclareMathOperator*{\colim}{colim}

\newcommand{\functorial}{\mathsf{ftl}}
\newcommand{\levelwise}{\mathsf{lw}}

\newcommand{\cofuniv}{\Phi}
\newcommand{\cofunivLW}{\Phi_{\levelwise}}
\newcommand{\cofunivConst}{\Phi}

\newcommand{\Gpd}{\mathsf{Gpd}}

\newcommand{\Modality}{\mathsf{Modality}}

\newcommand{\Set}{\mathsf{Set}}
\newcommand{\Psh}{\operatorname{PSh}}

% categories of maps
\newcommand{\Cof}{\mathsf{Cof}}
\newcommand{\TrivCof}{\mathsf{TrivCof}}
\newcommand{\Fib}{\mathsf{Fib}}
\newcommand{\TrivFib}{\mathsf{TrivFib}}
\newcommand{\GenCof}{\mathsf{I}}
\newcommand{\GenTrivCof}{\mathsf{J}}
\newcommand{\WEquiv}{\mathsf{W}}
\newcommand{\InjFib}{\mathsf{InjFib}}
\newcommand{\InjTrivFib}{\mathsf{InjTrivFib}}

\newcommand{\Alg}{\mathsf{Alg}}
\newcommand{\Coalg}{\mathsf{Coalg}}
\newcommand{\Map}{\mathsf{Map}}

\newcommand{\CofFun}{\mathsf{C}}
\newcommand{\TrivCofFun}{\mathsf{C_t}}
\newcommand{\FibFun}{\mathsf{F}}
\newcommand{\TrivFibFun}{\mathsf{F_t}}

\newcommand{\obj}{\operatorname{obj}}
\newcommand{\aug}{\mathsf{aug}}
\newcommand{\Endo}{\mathsf{Endo}}
\newcommand{\gen}{\mathsf{gen}}

\newcommand{\M}{\mathcal{M}}
\newcommand{\join}{\mathop{\star}}

\newcommand{\Prop}{\mathsf{Prop}}

\newcommand{\mycomment}[1]{}

\makeatletter % treat @ as a normal letter (catcode 11)
\newcommand{\varname}[1]{\begingroup\newmcodes@\mathit{#1}\endgroup} % \newmcodes resets \mathcode for `*-/:’'
\newcommand{\DeclareAbbrevation}[2]{\newcommand{#1}{\@ifnextchar{.}{#2}{#2.\@\xspace}}}
\makeatother  % treat @ not as a normal letter anymore (catcode 12)

%%% List environments
\newlist{conditions}{enumerate}{2}
\setlist[conditions,1]{label=\arabic*.,ref=\arabic*}
\setlist[conditions,2]{label=\arabic*.,ref=\arabic{conditionsi}.\arabic*}

\newlist{parts}{enumerate}{1}
\setlist[parts]{label=(\roman*)}

\newlist{options}{enumerate}{1}
\setlist[options]{label=\arabic*.,ref=\arabic*}

%%% Diagram environment
\newcounter{diagram}
\NewDocumentEnvironment{diagram}{oo}{
  \refstepcounter{diagram}
  \begin{center}
    \IfNoValueTF{#1}{
      \begin{tikzcd}
    }
    {
      \begin{tikzcd}[#1]
    }
}{
    \end{tikzcd}
    \IfNoValueF{#2}{
      \par\medskip
      Diagram~\thediagram: \textit{#2}
    }
  \end{center}
}
\NewDocumentEnvironment{diagram*}{o}{
  \begin{diagram}[#1]
}{
  \end{diagram}
}

%%% Pullback symbols
\newcommand{\pullback}[1]{\save*!/#1-1.2pc/#1:(-1,1)@^{|-}\restore}
\newcommand{\ulpullback}{\pullback{ul}}
\newcommand{\dlpullback}{\pullback{dl}}
\newcommand{\urpullback}{\pullback{ur}}
\newcommand{\drpullback}{\pullback{dr}}

%%% Delimiters
\DeclarePairedDelimiter\parens\lparen\rparen
\DeclarePairedDelimiter\floors\lfloor\rfloor
\DeclarePairedDelimiter\ceils\lceil\rceil
\DeclarePairedDelimiter\bracks\lbrack\rbrack
\DeclarePairedDelimiter\brackss\llbrack\rrbrack
\DeclarePairedDelimiter\verts\lvert\rvert
\DeclarePairedDelimiter\vertss\lVert\rVert
\DeclarePairedDelimiter\angles\langle\rangle
\DeclarePairedDelimiter\braces\lbrace\rbrace
\DeclarePairedDelimiterX\set[2]\lbrace\rbrace{#1 \mathrel{\delimsize\vert} #2}
\DeclarePairedDelimiter\ucorns\ulcorner\urcorner

%%% Abbreviations
\DeclareAbbrevation{\ie}{i.e}
\DeclareAbbrevation{\eg}{e.g}
\DeclareAbbrevation{\Eg}{E.g}
\DeclareAbbrevation{\cf}{cf}
\DeclareAbbrevation{\etc}{etc}
\DeclareAbbrevation{\resp}{resp}
\DeclareAbbrevation{\etal}{et al}
\DeclareAbbrevation{\ibid}{ibid}
\DeclareAbbrevation{\ca}{ca}
\DeclareAbbrevation{\vs}{vs}

%%% Cross-reference names
\crefname{conditionsi}{condition}{condition}
\Crefname{conditionsi}{Condition}{Condition}
\crefname{conditionsii}{condition}{condition}
\Crefname{conditionsii}{Condition}{Condition}
\crefname{diagram}{diagram}{diagrams}
\Crefname{diagram}{Diagram}{Diagrams}
\crefname{figure}{figure}{figures}
\Crefname{figure}{Figure}{Figures}
\crefname{formula}{formula}{formulas}
\Crefname{formula}{Formula}{Formulas}
\crefname{partsi}{part}{parts}
\Crefname{partsi}{Part}{Parts}
\crefname{optionsi}{option}{options}
\Crefname{optionsi}{Option}{Options}
\crefname{section}{Section}{Sections}
\Crefname{section}{Section}{Sections}
\crefname{square}{Square}{Squares}
\Crefname{square}{Square}{Squares}
\crefname{subsection}{Subsection}{Subsections}
\Crefname{subsection}{Subsection}{Subsections}
\crefname{subsubsection}{Subsection}{Subsections}
\Crefname{subsubsection}{Subsection}{Subsections}

%%% Theorem-like environments
\newtheorem{theorem}{Theorem}[section]
\newtheorem{corollary}[theorem]{Corollary}
\newtheorem{lemma}[theorem]{Lemma}
\newtheorem{proposition}[theorem]{Proposition}

\theoremstyle{definition}
\newtheorem{definition}[theorem]{Definition}

\theoremstyle{remark} % italic head, upright text, no extra space above or below
\newtheorem{remark}[theorem]{Remark}

%\AtAppendix{
%  \counterwithout{theorem}{section}
%}

\begin{document}

\title{Constructive Sheaf Models of Type Theory}

\author{Thierry Coquand, Fabian Ruch, and Christian Sattler}
\date{Computer Science Department, University of Gothenburg}
\maketitle

%\rightfooter{}

\section*{Introduction}

Despite being relatively recent, the notion of (pre\discretionary{-)}{}{)}sheaf model
has a rich and intricate history which mixes different intuitions
coming from topology, logic and algebra. Eilenberg and Zilber~\cite{EilenbergZ50}
used a presheaf model (simplicial sets) to represent
geometrical objects, and the intuition is {\em geometrical}:
we think of the objects $I,J,\dots$ of the base category as basic ``shapes'';
a presheaf $A$ is then given by a family of sets $A(I)$ of objects of each shape $I$,
which are related by the restriction maps $A(I) \rightarrow A(J)$.
A little later, but independently, Beth~\cite{Beth56} and Kripke~\cite{Kripke65}
used a sheaf and a presheaf model over trees, respectively,
to provide a formal semantics for intuitionistic logic. Their motivations
were logical, and the intuition is of a {\em temporal} nature instead: we think of the objects of the node of the tree
as ``stages of knowledge'' and of the ordering as ``increase in knowledge''.
Scott~\cite{Scott80} described a presheaf model of higher\hyp{}order logic and
pointed out the potential interest for the semantics of \textlambda\hyp{}calculus.
This was refined by Martin Hofmann~\cite{Hofmann97} who provided a presheaf
model of dependent type theory with universes. Hofmann's presheaf model was subsequently used in an
essential way in works on constructive semantics of type theory with
{\em univalent} universes~\cite{CCHM15,CHM18,OrtonP16}.

The generalization of such presheaf models of dependent type theory, and especially
of universes, to a {\em sheaf} model semantics
is however non\hyp{}trivial. The problem in generalising this semantics for universes
comes essentially from the fact that the collection of sheaves does not form a sheaf
in any natural way: if we are given locally sheaves that are compatible, one can patch
them together but not in a {\em unique} way, only unique {\em up to isomorphism}.
This problem was the motivation for the introduction of stacks and a more
subtle notion of patching of sheaves (\cf~\cite[Section~3.3]{EGAI}), and in general
patching of mathematical structures. The generalization of this to patching of
higher structures was the content of the first part of Joyal's letter to Grothendieck~\cite{Joyal84}.
One contribution of the present paper is to provide
a constructive version of this notion\footnote{Joyal's argument was using non\hyp{}constructive reasoning in simplicial sets
  and then Barr's theorem (see~\cite{Barr74}).
  The present paper can be developed directly in the constructive framework of CZF with universes
introduced by Aczel~\cite{Aczel98}.}
by describing a
sheaf model semantics of type theory with univalence~\cite{Voevodsky15,HoTT}.
This uses in a crucial way the fact that
we have a {\em constructive} interpretation of univalence as in~\cite{CCHM15,OrtonP16}, which can
be relativized to any presheaf model. The main point is then that the operation sending an object to its object of descent data (a compatible collection of elements of its restrictions) defines
a {\em left exact modality} (see~\cite{HoTT,Quirin16,RijkeSS17}), which can then be used
to build internally models of univalent type theory~\cite{Quirin16}.

This work opens the possibility of generalising works of sheaf models of intuitionistic
logic as in~\cite{TroelstraD88b} to sheaf models of univalent type theory. It extends the previous
work in~\cite{CoquandMR17} to a complete model of univalence, and has no restrictions for
representing (higher) data types.
We give only one application (independence of countable choice), but we expect for instance that results
such as in~\cite{MannaaC13} can be generalized as well, and that we can give a constructive
account of works such as in~\cite{Shulman18,Wellen17}.
The present semantics (in a preliminary version)
has already been used by Weaver and Licata~\cite{WeaverL20} for building a constructive model of directed univalence.

This paper is organized as follow. We first introduce the notion of
lex operation as an operation acting on types and families of types.
A descent data operation is then
a lex operation which defines a left exact modality \cite{HoTT,Quirin16,RijkeSS17}.
These two notions are formulated purely syntactically in the framework of type theory.
We show next how to instantiate these operations for cubical presheaves. In this setting,
we can understand the
notion of being modal for a descent data operation as a generalization of the sheaf condition,
where the compatibility requirements are expressed up to path equality instead of
being expressed as strict equalities. We then provide some examples and the application to
the unprovability of countable choice. In an appendix we explain how some
of our results about descent data operations can be generalized to accessible
left exact modalities.


%% The notion of non\hyp{}strict descent data is interesting even
%% for the trivial topology over a category: in this case we show that equivalences
%% in the resulting model coincide with pointwise equivalences.
%% A simple, but non\hyp{}trivial, example is the model of parameterized pointed types, where the base
%% category is the walking retract.

\section{Abstract notion of descent data}

In this article, we take terminology in type theory with potentially both strict and homotopical meaning to have the strict meaning by default.
For example, equality (denoted by the symbol ${=}$) refers to the strict equality (as opposed to identity or path types), isomorphisms refer to strictly invertible maps, and pullbacks refer to strict pullbacks.

We use the following notations.
We write $\Unit$ for the unit type and $\unit : \Unit$ for its unique element.
Given a type $A$ and a family $B$ of types over $A$, we write $\sum_A B$ for their sum type and $\prod_A B$ for their product type.
The pairing operation is denoted by $(a, b) : \sum_A B$ for $a : A$ and $b : B\,a$.
The projection maps are denoted by $\pi_1$ and $\pi_2$.
We write $\id_A$ for the identity function on $A$ and $g\circ f$ for the composition of $f \colon A \to B$ and $g \colon B \to C$.
If $B$ is a family of types over $A$ and $f\colon A'\rightarrow A$, we also write $B\circ f$ for the family of types over $A'$ obtained from $B$ by reindexing along $f$.

\subsection{Lex operation} \label{subsec:lex-operation}

The concept of lex operations is defined for a dependent type theory with only unit type, dependent sums, dependent product and universes.
In particular, path types are not needed.
Intuitively, a lex operation is an endofunctor on the category of types and functions (compatible with substitution) which preserves the unit type and dependent projections of sum types up to isomorphism.

A {\em lex operation}\footnote{The notion of lex operation appears implicitly in a natural way
when describing the rules of inductive data types~\cite{CoquandP88}. If we have a family
$\D_a$ of lex operations indexed over $a:A$, we can consider the inductive type $T$
with constructor $\msup:\prod_{a:A} (\D_aT\rightarrow T)$ and elimination rule
$\rec\,f:\prod_TP$ for $f:\prod_{a:A} \prod_{u:\D_aT} (\DD_aP\,u\rightarrow P(\msup_a\,u))$.
We can then write the computation rule
%
\begin{equation*}
\rec\,f\,(\msup_a\,u) = f_a\,u\,(\dd_a(\rec\,f)\,u)
\end{equation*}
%
For justifying the use of such inductive definitions, we need some ``accessibility'' assumption
on the functor $\D$, which will be satisfied in the examples.
In the special case where $\D_aX$ is $X^{Ba}$ for $B$ a family of over $A$ we recover
the $W$\hyp{}type $W_A B$.}
is given by an operation $\D$ on types and functions forming a functor:
we have $\D f:\D A\rightarrow \D B$ if $f:A\rightarrow B$ with $\D(g\circ f) = \D g\circ \D f$
and $\D(\id_A) = \id_{\D A}$.

The operation $\D$ should also preserve the unit type $\Unit$ up to isomorphism.
Specifically we have an element $\unitF$ in $\D\Unit$ and $x = \unitF$ if $x$ is in $\D\Unit$.

Furthermore, $\D$ should preserve dependent projections of sum types up to isomorphism.
We assert this by an operation on families of types: $\DD B$ is family of types over $\D A$
if $B$ is a family of types over $A$.
This should be natural in $A$
%\fxnote{stable under base change/natural in $A$
%(otherwise $\DD (B\circ f) = \DD B\circ \D f$ and naturality does not make sense)}
together with operations ensuring that $\D(\sum_AB)$ is naturally
isomorphic to $\sum_{\D A}{\DD B}$ over $\D A$.
Naturality in $A$ means $\DD (B\circ f) = \DD B\circ \D f:\D{A'}\rightarrow\UU$ for $f:A'\rightarrow A$.
%\fxnote{Note that we do not require $A$ to be in $\UU$ here (critical for our development).}
The natural isomorphism between $\D(\sum_AB)$ and $\sum_{\D A}{\DD B}$ over $\D A$ is given by
an operation $\dd s : \prod_{\D A}{\DD B}$ on sections $s:\prod_A B$
satisfying $\dd (s\circ f) = \dd s \circ \D f : \prod_{\D{A'}} \DD (B\circ f)$
%\fxnote{stable under base change (otherwise the below is not enough to get a natural isomorphism.)}
and a pairing operation $\pairF{u}{v} : \D(\sum_A B)$ for elements $u:\D A$ and $v:(\DD B)\,u$
satisfying
\begin{align*}
(\D\pi_1)\pairF{u}{v} &= u &
(\dd\pi_2)\pairF{u}{v} &= v &
\pairF {(\D\pi_1)\,w}{(\dd\pi_2)\,w} &= w
\end{align*}
where $w:\D(\sum_AB)$.

We also assume that universes reflect these operations.
This means we have $\D A:\UU$ if $A:\UU$ and $\DD B:\D A\rightarrow\UU$ if $B:A\rightarrow \UU$ and $A$ a type (crucially, $A$ need not be in $\UU$ here).

\medskip

The canonical example of a lex operation is exponentiation with a fixed type $R$
(assumed to be in all universes).
We define $\D A = A^R$, $(\DD B)u = \prod_{x:R}\,B(u\,x)$, and $(\dd s)u = \lambda_{x:R}\,s(u\,x)$.
The pairing is given by $\pairF{u}{v} = \lambda_{x:R}\,(u\,x,v\,x)$.

\begin{remark}
  Let $\UU$ be a universe.
  The action of the operation $\DD$ on $\UU$\hyp{}small families is uniquely determined by the universal case $L = \DD\,\id_\UU : \D \UU \to \UU$: we have (and can define) $\DD B = L\circ \D B$ with $\D B:\D A \rightarrow \D\UU$ for $B:A\rightarrow\UU$.
  This corresponds to the ``escaping'' function in Section 2.5 of \cite{SchreiberS14}.
  We can thus describe the action of $\DD$ on $\UU$\hyp{}small families and associated operations by requiring that $\D$ applied to the ``universal $\UU$\hyp{}small fibration'' $\sum_{X:\UU} X \to \UU$ is isomorphic to a ``$\UU$\hyp{}small fibration'' (a projection of a type in $\UU$), and that $\D$ preserves pullbacks of this map.
\end{remark}

\begin{proposition} \label{lex-operation-pointed}
  Any lex operation $\D$ is uniquely pointed.%
  \footnote{We owe this observation to Dan Licata.}
\end{proposition}

\begin{proof}
  We define $\eta_A\, a = (\D\epsilon_a)\,\unitF$ with $\epsilon_a = \lambda_{x:\Unit}\,a$.
  We then have for $f:A\rightarrow B$
  %
  \begin{equation*}
  (\D f)(\eta_A\, a) = (\D f\circ \D\epsilon_a)\,\unitF = (\D\epsilon_{f\,a})\,\unitF = \eta_B(f\,a)
  \end{equation*}
  %
  Note furthermore that this natural transformation $\eta_A$ is uniquely determined,
  since we should have $\eta_{\Unit}\,\unit = \unitF$ and so
  $\eta_A\, a = \eta_A(\epsilon_a\,\unit) = (\D\epsilon_a)(\eta_{\Unit}\,\unit) = (\D\epsilon_a)\,\unitF$.
\end{proof}

\begin{remark}\label{iso}
  For $T : \UU$, the map
  $(\DD \epsilon_T)\,\unitF\rightarrow \sum_{\D\Unit}\DD \epsilon_T\rightarrow
  \D(\sum_{\Unit}\epsilon_T)\rightarrow \D T$ is an isomorphism, as a composition
  of isomorphisms.
  As a consequence, the map
  %
  \begin{equation*}
  (\DD B)(\eta_A\, a)\longrightarrow \D(B\,a)\qquad v\longmapsto (\D\pi_2)\pairF{\unitF}{v}
  \end{equation*}
  %
  is an isomorphism for a type $A$, a family $B$ over $A$, and $a : A$.
\end{remark}

%% We assume furthermore that we have $(\DD B)(\eta_A\, a) = \D(B\,a)$ strictly for any $A:\UU$
%% and $B:A\rightarrow \UU$ and $a:A$. This further condition is not needed (\cf~\cite{Ruch20})
%% but will hold for the examples of lex operations we will analyse and it simplifies
%% some arguments.
%% \fxnote{%
%% This is a strange condition.
%% We have a natural isomorphism $\DD B(\eta_A(a)) \simeq \D(B(a))$.
%% If you postulate that both sides are equal, you should also require that the canonical map between them is the identity.
%% }

%% \fxnote{%
%% Suggestion.
%% State the natural isomorphism $\DD B(\eta_A(a)) \simeq \D(B(a))$ as a lemma.
%% Observe that it happens to be an identity in our examples (only true if non\hyp{}dependent functions are implemented as special case of dependent functions).
%% Rewrite all proofs to use this isomorphism instead of the equality (this does not complicate the proofs).
%% }

%% \Fxnote{%
%% The more primitive (equivalent) version is $\DD \epsilon_B (\unitF) \simeq \D B$ for any $B : \UU$ (an identity in our examples).
%% Note that $\unitF$ can be replaced by any element of any type, \ie $\DD \epsilon_B \simeq \epsilon_{\D B}.$
%% Perhaps rather state this as a lemma (proof is $\sum_{\D \Unit} \DD \epsilon_B \simeq \D(\sum_{x:\Unit} B)$ plus $\D$ preserves terminal types) and derive $\DD B(\eta_A(a)) \simeq \D(B(a))$ as a corollary.
%% }

%% %We define $L_\D:\D\UU\rightarrow\UU$ by $L_\D = \DD\id_{\UU}$. We then have
%% %$L_\D(\eta_{\UU}\,X) = (\DD\id_{\UU})(\eta_{\UU}\,X) = \D X$ for all $X$ in $\UU$.

\medskip

Note that $\eta_A\,a = \lambda_{x:R}\,a$ for the example $\D A = A^R$
where the lex operation is exponentiation.
The isomorphisms of \cref{iso} are identities in this example.%
\footnote{This assumes that function types are implemented via dependent products.}
In fact, this will happen for all the example lex operations we will consider in this article.

\begin{remark} \label{D-reflection}
Recall our assumption that universes reflect the operation $\D$ on types and the operation $\DD$ on families.
\Cref{iso} shows that only the reflection of $\DD$ is essential.
If $\D$ is not reflected, we can define an isomorphic operation $\D' A = \DD (\epsilon_A)\, \unitF$ on types that is reflected.
The remaining structure of $\D$ transports across the isomorphism to define a lex operation $\D'$.
\end{remark}

\begin{remark}
Let $\mathcal{E}$ be a category with families~\cite{Dybjer95} modelling our type theory.
A lex operation in $\mathcal{E}$ can be defined from a pseudomorphism of cwfs with universes~\cite{KaposiHS19} from $\mathcal{E}$ to itself that is pointed as an endofunctor.
When working externally with a model, this is a convenient way of constructing a lex operation in it.
Note that the given pointing is then reconstructed by \cref{lex-operation-pointed}.
\end{remark}

\begin{remark}
For readers familiar with Martin Hofmann's semantic methodology~\cite{Hofmann97}, we note a concise description of lex operations expressed internally in presheaves over the category of contexts.
The object $\Type$ of types has the structure of a cwf with universes (context extension is given by sum types).
Up to the discussion of \cref{D-reflection}, a lex operation is a pseudomorphism of cwfs with universes from $\Type$ to itself.
This definition can be written in the language of two\hyp{}level type theory~\cite{AnnenkovCKS17}.
\end{remark}

\subsection{\texorpdfstring{$\D$}{D}\hyp{}modal types}

The notion of lex operation is defined at the level of ``pure'' dependent type theory,
without assuming any notion of path types. In presence of path types, we automatically
have the following preservation property.

\begin{theorem}\label{lex}
  Let $\D$ be a lex operation.
  Then $\D$ preserves equivalences.
\end{theorem}

\begin{proof}
  Note that if $f_0$ and $f_1$ are path equal then so are $\D f_0$ and $\D f_1$ by path induction.
  It follows that if $f$ and $g$ are inverses, then so are $\D f$ and $\D g$.
\end{proof}

Avigad \etal~\cite{AvigadKL15} explain how to build a fibration category from a model of
dependent type theory. \Cref{lex} implies that any lex operation defines an
endomorphism of the associated fibration category. A lex operation preserves
all finite homotopy limits (\eg, contractible types, homotopy pullbacks, homotopy equalizers, homotopy fibers, \ldots).

\medskip

In presence of path types, we can also define the following important notion of modal types.

\begin{definition}
  A type $A$ is called $\D$\hyp{}modal if the unit map $\eta_A:A\rightarrow \D A$ is an equivalence.
\end{definition}

\begin{proposition}\label{sum}
  If $A$ is $\D$\hyp{}modal and $B$ is a family of types over $A$, then
  $B$ is a family of $\D$\hyp{}modal types over $A$ if, and only if, $T = \sum_AB$ is $\D$\hyp{}modal.
\end{proposition}

\begin{proof}
  Let $f$ be the map
  $T\rightarrow \sum_{\D A}{\DD B},\ (a,b)\mapsto (\eta_A\,a,\eta_{B\,a}\,b)$.
  Since $\eta_A$ is an equivalence,
  each map $\eta_{B\,a}$ is an equivalence if, and only if, the map $f$ is an equivalence~\cite{HoTT}.
  But $f$ is an equivalence if and only if $\eta_{T}$ is an equivalence.
\end{proof}

\subsection{Abstract notion of descent data}

%% The notion of lex operations was formulated for an arbitrary model of
%% type theory. For the next notion of descent data, we work in a model
%% of cubical type theory, where the path and identity types are obtained
%% from an interval type~\cite{CCHM15,OrtonP16,Boulier18}.

\begin{theorem}\label{main}
  The following conditions are equivalent, for a lex operation $\D$
  \begin{enumerate}
  \item $\D$ defines a modality as axiomatized in~\cite{Quirin16,RijkeSS17}
  \item the map $\D\eta_A$ is an equivalence, and $\D\eta_A$ and $\eta_{\D A}$
    are path equal
  \end{enumerate}
\end{theorem}

\begin{proof}
  The first condition implies the second using the results in~\cite{Quirin16,RijkeSS17}.

  Conversely, assume that the map $\D\eta_A$ is an equivalence, and $\D\eta_A$ and $\eta_{\D A}$
  are path equal. Then $\eta_{\D A}$ is an equivalence as well and each type $\D A$ is $\D$\hyp{}modal.
  \Cref{sum} shows that $\D$\hyp{}modal types are closed by dependent sum types.
  We thus only have to prove that the map
  %
  \begin{equation*}
  F:(\D A\rightarrow B)\longrightarrow (A\rightarrow B) \qquad f\longmapsto f\circ\eta_A
  \end{equation*}
  %
  is an equivalence if $B$ is $\D$\hyp{}modal~\cite{Quirin16}.

  Let $p_B$ be a map $\D B\rightarrow B$ such that $p_B\circ\eta_B$
  is path equal to $\id_B$. We define a map
  %
  \begin{equation*}
  G:(A\rightarrow B)\longrightarrow (\D A\rightarrow B)\qquad u\longmapsto p_B\circ \D u
  \end{equation*}
  %
  We then have $F(Gu) = p_B\circ \D u\circ\eta_A = p_B\circ\eta_B\circ u$ which is path equal to $u$
  and $G(Ff) = p_B\circ \D(f\circ \eta_A) = p_B\circ \D f\circ \D\eta_A$ which is path equal
  to $p_B\circ \D f\circ \eta_{\D A} = p_B\circ\eta_B\circ f$ which is path equal to $f$. Hence $G$
  is an inverse to $F$ and $F$ is an equivalence.
\end{proof}



\begin{definition}
  A {\em descent data operation} is a lex operation $\D$
  satisfying the equivalent conditions of \cref{main}.
\end{definition}

Note that the first condition of \cref{main} is a (homotopy)
proposition. The second condition is the one which will be convenient to verify
for the main examples.


%% \begin{definition}
%%   An {\em abstract notion of descent data} is a lex operation $\D$
%%   which is well\hyp{}pointed (see~\cite{Kelly80}) up to homotopy, \ie there is a path
%%   $\epsilon_A\,i$
%%   between $\epsilon_A\,0 = \D(\eta_A)$ and $\epsilon_A\,1 = \eta_{\D A}$, which is uniform along fibrations:
%%   $\D^2(\sigma)\circ (\epsilon_{T}\,i)$ is strictly equal to
%%   $(\epsilon_A\,i)\circ \D(\sigma)$ whenever $\sigma:T\rightarrow A$ is a fibration.
%% \end{definition}

%When $\D$ satisfy this well\hyp{}pointedness condition, we may say that $A$ is a {\em sheaf}
%instead of saying that $A$ is $\D$\hyp{}modal, and
We write $\GG_\D(A)$ for the type (proposition) expressing that $A$ is $\D$\hyp{}modal.

\subsection{Closure properties}

 Let $\D$ be a descent data operation. %\fxnote{Grammar? Descent operation?}

\begin{lemma}\label{suff}
  For the map $\eta_A:A\rightarrow \D A$ to be an equivalence, it is enough to have
  a patch function $p_A:\D A\rightarrow A$ such that $p_A\circ\eta_A$ is path equal to the identity of $A$.
\end{lemma}

\begin{proof}
  If $p_A$ is such a patch function, we have
  $\id_{\D A} = \D(\id_A) = \D(p_A\circ\eta_A) = \D p_A\circ \D\eta_A$ which is path equal
  to $\D p_A\circ\eta_{\D A} = \eta_A\circ p_A$. Hence $p_A$ is an inverse of $\eta_A$ and $\eta_A$
  is an equivalence.
\end{proof}

\begin{lemma}\label{DD}
For $B$ a family of types over $A$ and any $u:\D A$, the type $(\DD B)\,u$ is $\D$\hyp{}modal.
\end{lemma}

\begin{proof}
  Since $\D(\sum_AB)$ is $\D$\hyp{}modal, so is the isomorphic type
$\sum_{\D A}{\DD B}$.
Using \cref{sum}, we have that $(\DD B)\,u$
is a $\D$\hyp{}modal type for any $u:\D A$.
\end{proof}

\begin{proposition}\label{univ}
The type $\UU_\D = \sum_{\UU}\GG_\D$ is a $\D$\hyp{}modal type.
\end{proposition}

\begin{proof}
  Consider the diagram
\[
\xymatrix{
  \D\UU_\D \ar[dr]^{\DD\pi_1} \\
  \UU_\D \ar[u]^{\eta} \ar[r]_{\pi_1} & \UU \rlap{.}
}
\]
It commutes up to homotopy since $(\DD\pi_1)(\eta\, X)$ is isomorphic to
$\D(\pi_1\, X)$ by \cref{iso}, which
%\fxnote{Weaken this to $(\DD\pi_1)(\eta\, X) \simeq \D(\pi_1\, X)$ (and cite lemma) if previous suggestion is implemented}
is path equal to $\pi_1\, X$ for any $X:\UU_\D$ by univalence(!). Note also that $\pi_1$
is an embedding since $\GG_\D$ is a family of propositions.

Since $(\DD\pi_1)\, A$ is $\D$\hyp{}modal by \cref{DD} for any $A:\D\UU_\D$,
the map $\DD\pi_1:\D\UU_\D\rightarrow\UU$
factorizes through $\pi_1:\UU_\D\rightarrow\UU$ and the corresponding
map $\D\UU_\D\rightarrow\UU_\D$ is a left inverse of $\eta:\UU_\D\rightarrow \D\UU_\D$
since $\pi_1$ is an embedding. Hence $\UU_\D$ is $\D$\hyp{}modal by \cref{suff}.
\end{proof}

\subsection{Model associated to a descent data operation}
%\fxnote{Model associated to a family of \emph{descent operations}?}

We can now define an internal translation which provides a new model of
univalent type theory with higher inductive types
for any descent data operation $\D$, following the work in~\cite{Quirin16}. A
type $A,p$ of the new model
is a type $A$ {\em together with}
a proof $p$ that this type is $\D$\hyp{}modal, while an element of a pair $A,p$ is an element of $A$.

%% \begin{align*}
%%  \inc &: A \rightarrow MA \\
%%  \patch &: \D(MA) \rightarrow MA \\
%%  \leftinv &: \textstyle\prod_{x : MA} (\patch (\eta_{MA}\, x) =_{MA} x)
%% \end{align*}
%\end{minipage}
%\end{figure}

%% This is a purely syntactical process. We define,
%% where $pf$ denotes the proof that the first component is a sheaf (for
%% a type $A$, $[A].\pi_2$ will be a proof that $\sem{A}$ is a sheaf)

%% \begin{tabular}{lclllllcl}
%%  $[x]$ & = & $x$ & & & & $[\prod_{A} B]$ & = & $(\prod_{\sem{A}} \sem{B},pf)$ \\
%%  $[M\,N]$ & = & $[M]\,[N]$ &&&& $[\sum_{A} B]$ & = & $(\sum_{\sem{A}} \sem{B},pf)$ \\
%%  $[\lambda_{x:A}\,M]$ & = & $\lambda_{x:\sem{A}}\,[M]$ &&&& $[\Id_A(M,N)]$ & = & $(\Id_\sem{A}([M],[N]),pf)$ \\
%%  $[M.\pi_1]$ & = & $[M].\pi_1$ &&&& $[\UU]$ & = & $(\sum_{X:\UU} \GG(X),pf)$ \\
%%  $[M.\pi_2]$ & = & $[M].\pi_2$ &&&& \\
%%  $[M,N]$ & = & $[M],[N]$ &&&& $\sem{A}$ & = & $[A].\pi_1$
%% \end{tabular}

%We then have that if $x_1:A_1,\dots,x_n:A_n\vdash M:A$
%then $x_1:\sem{A_1},\dots, x_n:\sem{A_n}\vdash [M]:\sem{A}$.

%In order to get a model, what matters is to have an element of sheaf
%structure in each case, but the actual value of this element does not matter.

In order to interpret the type of natural numbers with the desired
computation rules (not covered in~\cite{Quirin16}),
we need to use the following higher inductive type:%
\footnote{To justify the use
  of such inductive definitions, we need some accessibility assumption on the functor
  $\D$ that will be satisfied in the examples.}
\begin{alignat*}{3}
  & \zero     & \quad & : & \quad & \Nat \\
  & \SUCC     & \quad & : & \quad & \Nat \rightarrow \Nat \\
  & \patch    & \quad & : & \quad & \D\,\Nat \rightarrow \Nat \\
  & \leftinv  & \quad & : & \quad & \textstyle\prod_{x : \Nat} \patch (\eta_\Nat\, x) =_{\Nat} x
\end{alignat*}

This is equivalent to the type $\D N$ where $N$ is the usual inductive
type with constructors $\zero$ and $\SUCC$, but the type $\D N$ does not satisfy
the required computation rules.

%% \begin{align}
%%   \zero    &: \Nat \\
%%   \SUCC    &: \Nat \rightarrow \Nat\\
%%   \patch   &: \D\Nat \rightarrow \Nat \\
%%   \leftinv &: \textstyle\prod_{x : \Nat} \patch (\eta_\Nat\, x) =_{\Nat} x
%% \end{align}
The same idea
%, which is mentioned in~\cite{RijkeSS17},
applies to the interpretation of other inductive types
such as the $W$\hyp{}type.
%% What is noteworthy is that this description of inductive types is actually
%% in some sense simpler in this higher setting than in the setting of ordinary $1$\hyp{}sheaves
%% (\cf the direct description of sheaves of $W$\hyp{}types in~\cite{BergM08}).
%% \fxnote{
%% But you will use the same description in the 1\hyp{}categorical setting (using a quotient inductive type).
%% So what is your point?
%% The reference~\cite{BergM08} is not really relevant; in sheaves, it can be read as constructing this quotient inductive type from an impredicative universe of propositions.
%% I guess that if you had a impredicative universe of propositions in the higher setting, you could also use it to justify this higher inductive types.
%% }

\medskip

It also works for higher inductive types. For instance
the suspension of a type $A$ will be defined as
\begin{alignat*}{3}
  & \north, \south    & \quad & : & \quad & T \\
%  & \south     & \quad & : & \quad & \susp {A} \\
  & \merid     & \quad & : & \quad & {A}\rightarrow \north =_{T} \south \\
  & \patch    & \quad & : & \quad & \D T \rightarrow T \\
  & \leftinv  & \quad & : & \quad & \textstyle\prod_{z : T} \patch (\eta_{T}\, z) =_{T} z
\end{alignat*}

Note that having $\D$ defined as a {\em strict} functor is essential for such definitions.

\subsection{Generalization to a family of descent data operations}\label{Suniv}
%\fxnote{Generalization to a family of \emph{descent operations}?}

More generally, if we have a {\em family} of descent data operations $\D_S$ indexed
by a given type $S:\Cov$, with corresponding maps $\eta^S_A:A\rightarrow \D_SA$,
we can consider $\GG_{\Cov}(A)$ to be the proposition $\prod_{S:\Cov} \GG_{\D_S}(A)$
and $\UU(\Cov)$ which is $\sum_{\UU}\GG_{\Cov}$. We use the slightly shorter notation $\UU_S$ to denote the previously defined type $\UU_{\D_S} = \sum_{\UU}\GG_{\D_S}$

We let the preorder $\D_1\leqslant \D_2$ on descent data operations mean that any $\D_1$\hyp{}modal
type is $\D_2$\hyp{}modal. We say that $\Cov$ is {\em filtered} if we have
$\exists_{S:\Cov}.\,\D_S\leqslant \D_{S_1}\wedge \D_S\leqslant \D_{S_2}$ for any
$S_1, S_2:\Cov$.%
\footnote{Existence is defined as the propositional truncation of the dependent sum type~\cite{HoTT}.}

\begin{theorem}\label{filter}
If $\Cov$ is filtered then $\UU(\Cov)$ satisfies $\GG_{\Cov}$.
\end{theorem}

\begin{proof}
%%   We have the following homotopy commuting diagram for any $\D_S\leqslant \D_{S_1}$ in $\Cov$
%%   \begin{displaymath}
%%     \xymatrix{
%%    \D_{S_1}\UU_{\Cov} \ar[rr]^{} \ar[ddr]|<<<<<<{\DD_{S_1}\pi_1} &   & \ar[ddl]|<<<<<<{\DD_{S_1}\pi_1} \D_{S_1}\UU_S \\
%%    \UU_{\Cov} \ar[u]^{\eta^{S_1}} \ar[rr]^{} \ar[dr]_{\pi_1} &    & \UU_S \ar[dl]^{\pi_1} \ar[u]_{\eta^{S_1}}\\
%%     & \UU &
%%    }
%% \end{displaymath}
  For any $\D_S\leqslant \D_{S_1}$ in $\Cov$, $\UU_S$ is $\D_S$\hyp{}modal by \cref{univ}
  and so $\D_{S_1}$\hyp{}modal,
  and hence $\eta^{S_1}:\UU_{S}\rightarrow \D_{S_1}\UU_{S}$ has an inverse.
  It follows that the map $\DD_{S_1}\pi_1:\D_{S_1}\UU_{\Cov}\rightarrow \UU$
  factorizes through $\UU_S\rightarrow\UU$ and
  hence that for any $A:\D_{S_1}\UU_{\Cov}$ the type $(\DD_{S_1}\pi_1)\,A$
  is $\D_S$\hyp{}modal.

  If $\Cov$ is filtered, this implies that the type
  $(\DD_{S_1}\pi_1)\,A$ is $\D_{S_2}$\hyp{}modal for any $S_2$ in $\Cov$. Hence
  the map $\DD_{S_1}\pi_1:\D_{S_1}\UU_{\Cov}\rightarrow \UU$
  factorizes through $\UU_{\Cov}\rightarrow \UU$
  and the corresponding map $\DD_{S_1}\UU_{\Cov}\rightarrow\UU_{\Cov}$
  is a left inverse of $\UU_{\Cov}\rightarrow \D_{S_1}\UU_{\Cov}$.
  Hence $\UU(\Cov)$ is $\D_{S_1}$\hyp{}modal for any $S_1$ in $\Cov$
  by \cref{suff}.
\end{proof}

This shows that for a family of descent data operations satisfying the hypothesis of \cref{filter},
we still get a model of univalent type theory (with higher inductive types), interpreting a type
as a type together with a proof that this type is modal for each descent data operation.
Of all type formers, only the universe has to deal with interaction between the elements of the given family of descent data operations.

\subsection{Example}\label{exp}

If $R$ is a {\em proposition}, then for the lex operation defined by $\D A = A^R$ the two maps
$\D\eta_A$ and $\eta_{\D A}$ are path equal equivalences and hence exponentiation
defines a descent data operation in that case.

 The next \lcnamecref{sec:presheaf-models} will define a new kind of descent data operation for any presheaf model.


\section{Cubical presheaf models}\label{sec:presheaf-models}

\subsection{Cubical models}\label{cubical-sets}

Cubical models are presheaf models of univalent type theory specified by two parameters, an \emph{interval object} $\II$ and a \emph{cofibration classifier} $\cofuniv$.
Formally, we say that a \emph{cubical model} is a presheaf category with the following structure, as in Orton and Pitts~\cite{OrtonP16}.
\begin{itemize}
\item
The interval object $\II$ is connected and has distinct points $0$ and $1$.
Exponentiation with $\II$ has a right adjoint.
We also assume that $\II$ has the structure of a bounded distributive lattice.%
\footnote{%
This assumption simplifies one of our arguments (\cref{key2}).
However, our results also apply to the Cartesian variation of cubical models of Angiuli \etal~\cite{AngiuliBCHHL}).
There, one removes this hypothesis and instead adds that the diagonal $\II \to \II \times \II$ is a cofibration.
}
\item
The universal cofibration $\top \colon 1 \to \cofuniv$ is a levelwise decidable inclusion.
In the internal language of presheaves, we will work with $\cofuniv$ as a universe of certain propositions and leave the decoding function (given by equality with $\top : \cofuniv$) implicit.
Isomorphic cofibrations are equal.%
\footnote{This assumption is not strictly speaking necessary, but simplifies the theory.}
The interval endpoint inclusions $0, 1 \colon 1 \to \II$ are cofibrations.
Cofibrations are closed under finite union (finite disjunction), composition (dependent conjunction), and universal quantification over $\II$.
\end{itemize}
It is then known, following the work in~\cite{CCHM15,OrtonP16,CHM18}, how to define a model of univalent type theory with higher inductive types.

\subsection{Presheaves in cubical models}\label{subsec:presheaves-in-cubical-sets}

For the remainder of this \lcnamecref{sec:presheaf-models}, we fix a cubical model given by presheaves over a small category $\BB$.
We refer to this as the \emph{base model} (for example, it can be cubical sets).
We write $I,J,K,\dots$ for the objects of $\BB$.
We have the interval object by $\II_{\BB}$ and the cofibration classifier $\cofuniv_{\BB}$.

%\fxnote{Investigate if $\CC$ can be an internal category in presheaves over $\BB$. Discuss in the conclusion.}

Let $\CC$ be another small category.
We write $X,Y,Z,\dots$ for its objects.
We describe possibilities for turning presheaves over $\CC \times \BB$ into a cubical model.
For the interval object $\II$, we simply take $\II(X,I) = \II_{\BB}(I)$.
For the cofibration classifier, we have two reasonable options:
\begin{options}
\item
The first example is simply to take $\cofunivConst(X,I) = \cofuniv_{\BB}(I)$.
\item
The second example is to define an element $\psi$ of $\cofunivLW(X,I)$ to be
a family $\psi_f$ in $\cofuniv_{\BB}(I)$ for $f \colon Y\rightarrow X$ such that
$\psi_f\leqslant \psi_g$ if furthermore $g \colon Z\rightarrow Y$.
We then define the restriction operation $\psi (f,l)$ to be the family
$\psi(f,l)_g = \psi_{f g} l$ for $f \colon Y\rightarrow X$ and $g \colon Z\rightarrow Y$.
\end{options}
The motivation for the second example is that if
$\cofuniv_{\BB}(I)$ is the collection of (decidable) sieves on $I$, then
$\cofunivLW(X,I)$ becomes the collection of (decidable) sieves on $(X,I)$.%
\footnote{Classically, this corresponds to having all monomorphisms as cofibrations.}

The interval object $\II$ and any of the choices $\cofunivConst$ and $\cofunivLW$ fit all the requirements listed in \cref{cubical-sets}.
This turns presheaves over $\CC \times \BB$ into a cubical model.
In particular, we get a model of univalent type theory (and higher inductive types).
We are going to analyse the model obtained using the choice $\cofunivConst$ for the cofibration classifier and then indicate how to adapt these results for $\cofunivLW$.

\medskip

In this model, a context $\Gamma$ is interpreted by a presheaf
over $\CC\times\BB$ so a family of sets $\Gamma(X,I)$ with suitable
restriction maps $\rho\mapsto \rho(f,l)$ with $f:Y\rightarrow X$ in $\CC$ and
$l:J\rightarrow I$ in $\BB$.
%We may simply write $\rho f$ instead of $\rho (f,\id_I)$
%and $\rho g$ instead of $\rho (\id_X,g)$.

A dependent type $A$ over $\Gamma$ is then given by a presheaf over the category
of elements of $\Gamma$: for any $\rho$ in $\Gamma(X,I)$ we have a set
$A\rho$ with suitable restriction maps $A\rho \rightarrow A\rho(f,l)$ denoted by $u\mapsto u(f,l)$
together with a filling operation (see~\cite{CCHM15,OrtonP16}).
We write $\Type(\Gamma)$ for the collection of all types with a composition operation
over $\Gamma$.
The set $\Elem(\Gamma,A)$ is then the set of sections: a family $a\rho$
in $A\rho$ such that $(a\rho)(f,l) = a(\rho(f,l))$ for any $\rho$
in $\Gamma(X,I)$ and $f,l$ map of codomain $X,I$.

%% If $J$ is an object of $\BB$, we write $J^+$ for $J\times {\bf I}$ and
%% $f^+:K^+\rightarrow J^+$ for $f\times {\bf I}$ if $f:K\rightarrow J$.
%% We have two maps $\epsilon_0, \epsilon_1:J\rightarrow J^+$ and a projection
%% $\pp:J^+\rightarrow J$. A {\em composition operation} (see~\cite{CCHM15,OrtonP16}) for
%% $A$ in $\Type(\Gamma)$ is an operation $c_A$ which takes $\rho$ in
%% $\Gamma(X|J^+)$ and $\psi$ in $\cofuniv(X|J)$ and
%% $u_0$ in $A\rho\epsilon_0$ and
%% a family $u_{f|g}$ in $A\rho(f|g)$ for $\psi (f|g) = 1$ for
%% $f:Y\rightarrow X$ and $g:K\rightarrow J^+$ such that
%% $u_{f|g\epsilon_0} = u_0(f|g)$ if $f:Y\rightarrow X$
%% and $g:K\rightarrow J$ and $\psi(f|g) = 1$ and produces
%% an element $u_1 = c_A(\psi\rightarrow u)\,u_0$ in
%% $A\rho\epsilon_1$ such that
%% $u_{f|g\epsilon_1} = u_1(f|g)$ if $f:Y\rightarrow X$
%% and $g:K\rightarrow J$ and $\psi(f|g) = 1$. Furthermore
%% we should have the uniformity condition
%% $(c_A(\psi\rightarrow u)\,u_0)(f|g) = c_A(\psi(f|g)\rightarrow u(f|g^+))\,u_0(f|g)$.

Given a constructive Grothendieck universe $U$ (see~\cite{Aczel98}) containing $\BB$ and $\CC$,
we write $\Type_{U}(\Gamma)$ for the set of $U$\hyp{}types, such that each set
$A\rho$ is in $U$. The presheaf $\Type_U$ is then represented by a fibrant
type $\UU$, which is univalent~\cite{CCHM15}.

\subsection{Internal language description}\label{subsec:internal-language-description}

This was an external description of the presheaf model.
It is also possible to describe this model
using
the internal logic of the presheaf topos over $\CC\times\BB$ as in~\cite{OrtonP16,CHM18}
but also using
the internal logic of the presheaf topos over $\BB$.
We will use both descriptions.

In the internal logic of the presheaf topos over $\BB$,
a context of the presheaf model over $\CC$ is interpreted as a family of ``spaces'' $\Gamma(X)$
with restriction maps $\rho\mapsto\rho f$ for $f:Y\rightarrow X$.
(Each space $\Gamma(X)$ is itself a presheaf over $\BB$ with $\Gamma(X)(I) = \Gamma(X,I)$.)
A dependent type $A$ over $\Gamma$ is given by a family of spaces $A\rho$
for $\rho$ in $\Gamma(X)$ with restriction maps $u\mapsto uf$. The presheaf
$\cofunivConst$ of cofibration is the constant presheaf $\cofunivConst(X) = \cofuniv_{\BB}$.
The interval $\II$ is the constant interval $\II(X) = \II_{\BB}$.

It will be convenient to introduce the following notation: if
$\gamma$ is an element of $\Gamma(X)^{\II}$ and
$f:Y\rightarrow X$, we write $\gamma f^+$ in $\Gamma(Y)^{\II}$
for $\lambda_i\,\gamma(i)f$. Similarly if $u(i)$ is a section
in $A\gamma(i)$ we write $uf^+$ for $\lambda_i\,u(i)f$.


% We write $I$ the identity map $I\rightarrow I$.

A {\em filling operation}
(see~\cite{OrtonP16,CHM18}) for $A$ is given by an operation $c_A$ which takes as
argument $\gamma$ in $\Gamma(X)^{\II_{\BB}}$ and $\psi$ in $\cofunivConst(X) = \cofuniv_{\BB}$
and a family of elements $u(i)$ in $A\gamma(i)f$ on the extent
$\psi\vee i = 0$. (There is a dual operation with $i=1$ instead.)
It produces an element $c_A(X,\gamma,\psi,u)(i)$ in $A\gamma(i)$ such that
\begin{conditions}
\item $c_A(X,\gamma,\psi,u)(i) = u(i)$ on $\psi\vee i=0$,
\item $c_A(X,\gamma,\psi,u)(i)f = c_A(Y,\gamma f^+,\psi, u f^+)(i)$ for $f:Y\rightarrow X$.
\end{conditions}
Given such an operation, we also call $A$ \emph{fibrant} (note this is structure rather than property).

If $A$ is a type over $\Gamma$, we get a family of dependent types $A(X)$
over $\Gamma(X)$, each of them having a filling operation, but furthermore
these filling operations commute with the restriction maps.

Similarly an {\em extension operation} for $A$, witnessing that $A$ is contractible
(see~\cite{CCHM15}), is given by an operation $e_A$ which takes as argument $\rho$
in $\Gamma(X)$ and a partial element $u$ on the extent $\psi$ and produces
an element $e_A(X,\rho,\psi,u)$ in $A\rho$ such that
\begin{conditions}
\item $e_A(X,\rho,\psi,u) = u$ on $\psi$,
\item $e_A(X,\rho,\psi,u)f = e_A(Y,\rho f,\psi, uf)$ for $f:Y\rightarrow X$.
\end{conditions}
Given such an operation, we also call $A$ \emph{trivially fibrant} (again, this is structure rather than property).

If $A$ is contractible, each $A(X)$ is a contractible family of types over $\Gamma(X)$.
But conversely, it may be that each $A(X)$ has an extension operation $e_A(X)$ which
does not commute with restriction (see \cref{subsec:examples}).
Similarly, a map $\sigma:A\rightarrow B$
which is an equivalence defines a family of equivalences $\sigma_X:A(X)\rightarrow B(X)$
but it may be that each map $\sigma_X$ is an equivalence, without $\sigma$ being
an equivalence.

\begin{remark} \label{presheaf-model-comparison}
We have a canonical map from $\cofunivConst$ to $\cofunivLW$ sending $\psi : \cofunivConst(X)$ to the constant family on $\psi$.
This map commutes (up to isomorphism) with the decoding to propositions.
It follows that there is a natural map from extension operations for $\cofunivConst$ to extension operations for $\cofunivLW$, and the same holds for filling operations.
It follows that a (contractible) type for the cubical presheaf model for $\cofunivConst$ is naturally also a (contractible) type for the cubical presheaf model for $\cofunivLW$.
\end{remark}

\begin{remark} \label{presheaf-model-groupoid}
Let $\CC$ be a groupoid.
Then for $\psi : \cofunivLW(X)$ and $f : Y \to X$, we have $\psi_{\id_X} \leq \psi_f \leq \psi_{f f^{-1}} = \psi_{\id_X}.$
It follows that $\psi$ is the constant family on $\psi_{\id_X}$.
Thus, the map $\cofunivConst \to \cofunivLW$ from \cref{presheaf-model-comparison} is invertible.
It follows that the cubical presheaf models for $\cofunivConst$ and $\cofunivLW$ are the same.
We thank Emily Riehl for this observation.
\end{remark}

\subsection{Examples} \label{subsec:examples}

Let $\BB$ be a concrete cube category, for instance the Cartesian~\cite{AngiuliBCHHL}, distributive lattice, or de Morgan one~\cite{OrtonP16,CCHM15}.
Then we have a nerve functor from groupoids to cubical sets in the sense of presheaves over $\BB$.
In this way, we can see any groupoid as a cubical set with a canonical filling operation.

For the first example, let $\CC$ be the group $\ints/(2)$.
Let $\tau$ be the non\hyp{}trivial element of this group.
A context is a space with an involutive action $\rho\mapsto\rho\tau$.
A dependent type $A$ over $\Gamma$ has also an involutive action $A\rho\rightarrow A\rho\tau$ denoted by $u\mapsto u\tau$ with a filling operation which is equivariant, meaning $c_A(\gamma,\psi,u)(i)\tau = c_A(\gamma \tau^+,\psi, u\tau^+)(i)$.
Let $A$ be the groupoid with two isomorphic objects swapped by $\tau$.
Then $A$ is pointwise contractible, but is not contractible in the presheaf model, since it has no global point.
Another way to describe this example is that the unique map $A\rightarrow \Unit$ is a pointwise equivalence, but is not an equivalence.

For the second example, let $\CC$ be the poset on objects $\bot, 0, 1$ with $\bot < 0$ and $\bot < 1$.
We define a global type $A$ as follows.
We take $A(0)$ and $A(1)$ to consist of a single object $a_0$ and $a_1$, respectively.
We take $A(\bot)$ to consist of an isomorphism between the restrictions of $a_0$ and $a_1$.
Then $A$ is levelwise contractible (\ie, $A(\bot), A(0), A(1)$ are contractible), but $A$ is not contractible since it has no global point.

We note that the second example is fixed by working with the cofibration classifier $\cofunivLW$.
However, as explained by \cref{presheaf-model-groupoid}, this does not apply to the first example.

\section{Homotopy descent data}

\subsection{A lex operation}
\label{sec:a-lex-operation}

In this \lcnamecref{sec:a-lex-operation}, we work in the internal language of the presheaf topos
over $\BB$.
We first define a lex operation on presheaf types, and then show that
this lex operation extends to types with a filling operation.

For any $A$ presheaf over $\Gamma$ we define $\E A$ presheaf over $\Gamma$.
An element $u$ of $(\E A)\rho$, for $\rho$ in $\Gamma(X)$ is given
by a family of elements $u(f)$ in $A\rho f$ for $f:Y\rightarrow X$.
We define the restriction $uf$ in $(\E A)\rho f$ by $uf(g) = u(fg)$
if $f:Y\rightarrow X$ and $g:Z\rightarrow Y$.

If $B$ is presheaf over $\Gamma.A$, we define $\EE(B)$ presheaf over  $\Gamma.{\E A}$. If
$\rho$ is in $\Gamma(X)$ and $u$ is in $(\E A)\rho$, then $\EE(B)(\rho,u)$ is the space of families
$v(f)$ in $B(\rho f,u(f))$.

We define a natural transformation $\alpha:A\rightarrow \E A$ by $(\alpha a)(f) = af$.

Next, we extend the action of $\E$ to types with a filling operation. Actually,
we define a filling operation $\E(c_A)$ on $\E A$ assuming only that $c_A$
is a {\em pointwise} filling operation on $A$.

\begin{proposition}\label{Efill}
  We can define a filling operation $\E(c_A)$ on $\E A$ if $c_A$ is a pointwise
  filling operation on $A$, in a way which commutes with substitution.
\end{proposition}

\begin{proof}
  We assume that $A$ has a pointwise filling operation $c_A(X)$.
  We define then, for $f:Y\rightarrow X$
  %
  \begin{equation*}
  \E(c_A)(X,\gamma,\psi,u)(i)(f) = c_A(Y)( \gamma f^+,\psi,u f^+)(i)
  \end{equation*}
  %
  We can then check for $f:Y\rightarrow X$ and $g:Z\rightarrow Y$
  %
  \begin{equation*}
  c_{\E A}(X,\gamma,\psi,u)(i)f(g) =
  c_A(Z)( \gamma (f g)^+,\psi,u (f g)^+)(i) =
  c_{\E A}(Y, \gamma f^+,\psi, u f^+)(i)(g)
  \end{equation*}
  %
  and hence $c_{\E A}$ is natural in $X$.

  We can also define $\unitF$ in $\E \Unit$ by $\unitF(f) = \unit$
  and $\pairF{u}{v}:\E(\sum_A B)\rho$
  by $\pairF{u}{v}(f) = (u(f),v(f))$ for $u$ in $(\E A)\rho$ and
  $v$ in $(\EE B)(\rho,u)$, and check that all conditions for a lex operations are satisfied.

  Any universe $\UU$ reflects the operations $\E$ and $\EE$ since the Grothendieck universe $U$ used to construct $\UU$ was assumed to contain $\CC$.
\end{proof}

\begin{proposition}\label{Econtr}
  If $A$ is pointwise contractible then $\E A$ is contractible.
\end{proposition}
\begin{proof}
  We assume that $A$ has a pointwise extension operation $e_A(X)$.
  We define then, for $f:Y\rightarrow X$
  %
  \begin{equation*}
  e_{\E A}(X,\rho,\psi,u)(f) = e_A(Y)(\rho f,\psi,uf)
  \end{equation*}
  %
  We can then check for $f:Y\rightarrow X$ and $g:Z\rightarrow Y$
  %
  \begin{equation*}
  e_{\E A}(X,\rho,\psi,u)f(g) =
  e_A(Z)(\rho f g,\psi,u f g) =
  e_{\E A}(Y,\rho f,\psi,u f)(g)
  \end{equation*}
  %
  and hence $e_{\E A}$ is an extension operation for $\E A$ natural in $X$.
\end{proof}

% If $u$ is in $(\E^mA)(X)$ and $f_1,\dots,f_m$ composable chain of arrows from $X$
% in $\CC$, then the element
% %
% \begin{equation*}
% (\E^k(\alpha)\,u)(f_1,\dots,f_k,f_{k+1},\dots,f_m)
% \end{equation*}
% %
% is equal to $u(f_1,\dots,f_kf_{k+1},\dots,f_m)$ for $k<m$, while the element (for $m>0$)
% %
% \begin{equation*}
% (\E^n(\alpha)\,u)(f_1,\dots,f_{m-1},f_m)
% \end{equation*}
% %
% is equal to $u(f_1,\dots,f_{m-1})f_m$.

In general, $\E$ may not be a descent data operation, since $\E A$ does not need to be $\E$\hyp{}modal.
The next \lcnamecref{sec:homotopy-descent-data} will use the lex operation $\E$ to define a descent data operation.
%is not a descent data operation: we may not have a path between
%$\E(\alpha_A)$ and $\alpha_{\E A}$.

\subsection{Homotopy descent data}
\label{sec:homotopy-descent-data}

In this \lcnamecref{sec:homotopy-descent-data}, unless explicitly stated,
we work in the internal language of the presheaf model over $\CC\times\BB$.
Starting from the lex operation $\E$, we define a new lex operation
$\D$.
As before, we first define $\D$ on presheaves, and then show that it extends
to a lex operation on presheaves with a filling operation.
On presheaves with a filling operation,
$\D$ will be  a descent data operation.

We let $\PP_n$ be the subpresheaf
of $\II^{n+1}$ of elements $(i_0,i_1,\dots,i_n)$ satisfying $i_0=1\vee\dots\vee i_n=1$.

\medskip

Let $s_k:\II^{n+1}\rightarrow \II^n$ be the map which omits the $k$th component,
for $k = 0,\dots, n$. Note that $s_k\vec{i}$ is in $\PP_{n-1}$
if $\vec{i}$ is in $\PP_n$ and $i_k=0$.

%% We define $\delta(i_1,\dots,i_n)$ in $\cofuniv$ as the disjunction of
%% $0 = i_1$ and all $i_k = i_{k+1}$ and $i_n = 1$.
%% We define then the partial element $\diag(i_1,\dots,i_n)$
%% of extent $\delta(i_1,\dots,i_n)$ as the element of $\II^{n-1}$
%% obtained from $i_1,\dots,i_n$ by omitting $i_1$ on $0=i_1$ and omitting $i_{k+1}$ on $i_k = i_{k+1}$
%% and omitting $i_n$ on $i_n =1$.

\begin{definition}
  An element of $\D A$ is given by a family %$u()$, $u(i)$, $u(i,j)$,\ldots with
  $u(\vec{i})$ in $\E^{n+1}A$ defined on $\PP_n$
  and satisfying the {\em compatibility conditions}\footnote{It is suggestive to think of the elements
    of $\D A$ as {\em choice sequences} \cite{TroelstraD88b} extended in a {\em spatial} rather than
    {\em temporal} dimension.}
%  $u(0) = \alpha u()$, $u(1) = \E(\alpha) u()$, $u(0,j) = \alpha u(j)$, $u(i,i) = \E(\alpha) u(i)$,
%  $u(i,1) = \E^2(\alpha) u(i)$,\ldots and more generally
$u(\vec{i}) = \E^k(\alpha) u(s_k\vec{i})$ on $i_k=0$.
\end{definition}

For instance we have
\begin{align*}
u(0,i_1,i_2) &= \alpha u(i_1,i_2) &
u(i_0,0,i_2) &= \E(\alpha) u(i_0,i_2) &
u(i_0,i_1,0) &= \E^2(\alpha) u(i_0,i_1)
\end{align*}

We have an element $u(\vec{1})$ in each $\E^{n+1}A$.
We have a path $u(1,i)$ between $\alpha\,u(1)$ and $u(1,1)$
and a path $u(i,1)$ between $\E(\alpha)\,u(1)$ and $u(1,1)$ in $\E^2A$.
But, in general, we need further higher coherence conditions.

\medskip

We define $\eta_A : A\rightarrow \D A$ by $(\eta_A\, a)(i_0,i_1,\dots,i_n) = \alpha^{n+1} a$.

If $A$ is a family of types over $\Gamma$ we define $\D A$ family of types
over $\Gamma$ by $(\D A)\rho = \D(A\rho)$.

\begin{proposition}\label{Dfill}
  If $A$ is a family of types with a pointwise filling operation,
  then $\D A$ has a filling operation.
\end{proposition}

\begin{proof}
  We use that each $\E^{n+1}A$ has a (uniform) filling operation by \cref{Efill}
  hence is a family of types in the model over $\CC\times\BB$.
  We assume given $\gamma$ in $\Gamma^{\II}$ and $\psi$ in $\cofunivConst$ and a partial
  element $u_j$ in $(\D A)\gamma(j)$ defined over $\psi\vee j=0$. We explain
  how to define a total extension $v_j$ in $(\D A)\gamma(j)$.
  For this we define $v_j(\vec{i})$ in $\E^{n+1}A$  by induction on $n$.
  Since $\E^{n+1}A$ has a filling operation, we apply this filling operation
  to the partial element equal to $u_j(\vec{i})$ on $\psi\vee j=0$
  and equal to $\E^k(\alpha)\, {v_j}(s_k(\vec{i}))$ if $i_k = 0$.
\end{proof}


\begin{corollary}\label{Dlex}
  $\D$ defines a lex operation.
\qed
\end{corollary}

A similar argument as the one for \cref{Dfill}
using \cref{Econtr} instead proves the following.

\begin{proposition}\label{Dcontr} \leavevmode
\begin{parts}
\item
If $A$ is a pointwise contractible family of types over $\Gamma$, then $\D A$ is contractible.
\item
If $B$ is a pointwise contractible family of types over a family of types $A$ over $\Gamma$, then $\DD B$ is contractible over $\D A$.
\qed
\end{parts}
\end{proposition}

\begin{corollary}\label{Dequiv}
Let $\sigma : A \to B$ be map between fibrant families of types over $\Gamma$.
If $\sigma$ is pointwise an equivalence, then $\D\sigma$ is an equivalence.
\end{corollary}

\begin{proof}
  The fiber $\hFiber(\sigma)$ defines a pointwise contractible family
  of types over $B$. Hence $\DD \hFiber(\sigma)$ is contractible over $\D B$.
  Since $\D$ is a lex operation,  $\hFiber(\D\sigma)$ is contractible over $\D B$
  and $\D\sigma$ is an equivalence.
\end{proof}

\begin{proposition}\label{key1}
Let $A$ be a fibrant family of types over $\Gamma$.
Then $\eta_A$ is pointwise an equivalence and $\D\eta_A$ is an equivalence.
\end{proposition}

\begin{proof}
  For this \lcnamecref{key1}, we work in the presheaf model over $\BB$.
  If $\vec{f}$ is a composable chain of arrows we write $\mpF{\vec{f}}$ for its composition.

Let $A$ be a type over $\Gamma$. For $\rho$ in $\Gamma(X)$, an element $u$
  of $(\D A)\rho$ is a family of elements $u(\vec{i})(\vec{f})$ in $A\rho \mpF{\vec{f}}$
  satisfying   the compatibility conditions.
  For $a$ in $A\rho$ the element $\eta_A\, a$ is the family
  of element
  %
  \begin{equation*}
  (\eta_A\, a)(\vec{i})(\vec{f}) = a\mpF{\vec{f}}
  \end{equation*}
  %
  We define an inverse $G : DA(X) \to A(X)$ of $\eta_A(X)$ by taking $Gu$ to be the element
  $u(1)(\id_X)$. We then have $G(\eta_A\, a) = a$. The element
  $\eta_A\, (G\,u)$ satisfies
  %
  \begin{equation*}
  (\eta_A\, (G\,u))(\vec{i})(\vec{f}) = (G\,u)\mpF{\vec{f}} = u(1)(\id)\mpF{\vec{f}}
  = u(1,\vec{0})(\id,\vec{f})
  \end{equation*}
  %
  Define the element $\tilde{u}$ in $(\D A)\rho$ by
  $\tilde{u}(\vec{i})(\vec{f}) = u(1,\vec{i})(\id,\vec{f})$.
  We can define a homotopy
  %
  \begin{equation*}
  u_k(\vec{i})(\vec{f}) = u(1,k\wedge \vec{i})(\id,\vec{f})
  \end{equation*}
  %
  between $\eta_A\,(G\,u)$ and $\tilde{u}$ and we can define a homotopy
  %
  \begin{equation*}
  v_k(\vec{i})(\vec{f}) = u(k,\vec{i})(\id,\vec{f})
  \end{equation*}
  %
  between $u$ and $\tilde{u}$.%
  \footnote{At this point that we use that the object {\bf I} in $\BB$ has lattice
  operations but one could however instead define a homotopy in a more complex way by induction on
  the dimension for Cartesian cubes. The same remark applies for the proof of the next
  \lcnamecref{key2}.}
  By composition, there is a path between $u$ and $\eta_A\, (G\,u)$
  and $G$ is an inverse of $\eta_A(X)$.%

  This shows that $\eta_A$ is pointwise an equivalence.
  Then $D \eta_A$ is an equivalence by \cref{Dequiv}.
\end{proof}

One way to understand the definition of $\D$ from $\E$ is the following.
Being a pointed endofunctor, $\E$ defines a cosemisimplicial
diagram starting from $\E A$, and $\D A$ is a strict way to realize
the homotopy limit of this diagram using a $\PP$\hyp{}weighted limit.
We can think of $\PP$ as a cofibrant resolution of the constant diagram on $1$.
A remark is that $\E$, and hence each $\E^l$, preserves the $\PP$\hyp{}weighted limit defining $\D$.
In particular, an element of $\E^l(\D A)$ is determined
by a family %$u()$, $u(i)$, $u(i,j)$,\ldots with
$u(\vec{i})$ in $\E^{l+n+1}A$ satisfying
%  $u(0) = \alpha u()$, $u(1) = \E(\alpha) u()$, $u(0,j) = \alpha u(j)$, $u(i,i) = \E(\alpha) u(i)$,
%  $u(i,1) = \E^2(\alpha) u(i)$,\ldots and more generally
$u(\vec{i}) = \E^{l+k}(\alpha)\, u(s_k\vec{i})$ on $i_k = 0$.

\begin{proposition}\label{key2}
  Let $A$ be a fibrant family of types over $\Gamma$.
  We can build a path between $\eta_{\D A}$ and $\D\eta_A$.
\end{proposition}

\begin{proof}
  An element of $(\D^2A)\rho$ is given by a family $v(\vec{i})(\vec{j})$
  in $\E^{n+m+2}A$
  %$v(i_1,\dots,i_n)(j_1,\dots,j_m)$
  satisfying the conditions
\begin{enumerate}
\item $v(\vec{i})(\vec{j}) = \E^k(\alpha)\, v(s_k\vec{i})(\vec{j})$ on $i_k=0$
\item $v(\vec{i})(\vec{j}) = \E^{n+1+l}(\alpha)\, v(\vec{i})(s_l\vec{j})$ on $j_l=0$
\end{enumerate}
Given $u$ in $(\D A)\rho$ we define an element $\tilde{u}$ in $(\D^2A)\rho$
by $\tilde{u}(\vec{i})(\vec{j}) = u(\vec{i},\vec{j})$.

We compute, for $u$ in $(\D A)\rho$
%
\begin{equation*}
(\eta_{\D A}\, u)(\vec{i})(\vec{j}) = \alpha^{n+1}\, u(\vec{j}) = {u}(\vec{0},\vec{j})
\end{equation*}
%
and we have a homotopy connecting this map to $\tilde{u}$ by defining
%
\begin{equation*}
v_k(\vec{i})(\vec{j}) = u(\vec{i}\wedge k,\vec{j}).
\end{equation*}
%
We also have
%
\begin{equation*}
((\D\eta_A)\, u)(\vec{i})(\vec{j}) =
\E^{n+1}(\alpha^{m+1})\,u(\vec{i}) = {u}(\vec{i},\vec{0})
\end{equation*}
%
and we have a homotopy connecting this map to $\tilde{u}$ by defining
%
\begin{equation*}
w_k(\vec{i})(\vec{j}) = u(\vec{i},k\wedge \vec{j})
\end{equation*}
%
By composition, we have a path between $\D\eta_A$ and $\eta_{\D A}$.
\end{proof}

\begin{corollary}
$\D$ defines a descent data operation.
\end{corollary}

\begin{proof}
  By \cref{key1,key2}.
\end{proof}

Note that a direct consequence of \cref{Dequiv} is the following
strictification result.

\begin{theorem}\label{strict1}
  Let $A$ and $B$ be fibrant families of types over $\Gamma$ that are $\D$\hyp{}modal.
  Then any pointwise equivalence $\sigma : A \to B$ is an equivalence.
\qed
\end{theorem}

Let us note the following consequence of \cref{key1}.

\begin{corollary} \label{all-modal-via-pointwise-contractible}
The following conditions are equivalent:
\begin{conditions}
\item \label{all-modal-via-pointwise-contractible:modal}
all fibrant families of types are $D$-modal,
\item \label{all-modal-via-pointwise-contractible:euqivalence}
all pointwise equivalences between fibrant families of types are equivalences,
\item \label{all-modal-via-pointwise-contractible:contractible}
all fibrant families of types that are pointwise contractible are contractible.
\end{conditions}
\end{corollary}

\begin{proof}
The direction from \cref{all-modal-via-pointwise-contractible:modal} to \cref{all-modal-via-pointwise-contractible:euqivalence} is \cref{strict1}.
In the reverse direction, given a fibrant family of types $A$, recall that $\eta_A$ is a pointwise equivalence by \cref{key1}.
Then $\eta_A$ is an equivalence and hence $D$-modal.
\Cref{all-modal-via-pointwise-contractible:contractible} is a special case of \cref{all-modal-via-pointwise-contractible:euqivalence}.
The reverse direction holds since a (pointwise) equivalence can be described as a map with (pointwise) contractible fibers.
\end{proof}

The way from which we get $\D$ from $\E$ can also be applied to the
lex operation $\E A = A^R$, where $R$ is an arbitrary type.
%% An element of
%% $\D A$ will be a family of maps $u(\vec{i}):R^{n+1}\rightarrow A$
%% satisfying $u(0)(r_1,r_2) = u()(r_2)$, $u(1)(r_1,r_2) = u()(r_1)$
%% and
%% $u(0,j)(r_1,r_2,r_3) = u(j)(r_2,r_3)$ and $u(i,i)(r_1,r_2,r_3) = u(i)(r_1,r_3)$
%% and $u(i,1)(r_1,r_2,r_3) = u(i)(r_1,r_2)$ and so on.
This amounts to give a map which is {\em coherently constant} as defined by Kraus~\cite{Kraus15}
and so a map $\norm{R}\rightarrow A$ from the propositional truncation
of $R$ to $A$~\cite{Kraus15}.

Our development actually provides a way to recover
this result in the cubical setting.
Indeed, an element of $\D A$ is a sequence of elements $u(\vec{i})(\vec{x})$
in $A$ for $\vec{i}$ in $\PP_n$ and $\vec{x}$ in $R^{n+1}$ with
$u(\vec{i})(\vec{x}) = u(s_k\vec{i})(s_k\vec{x})$ on $i_k = 0$. Given an element
$x$ in $R$, we can build a left inverse $p_A$ of $\eta_A:A\rightarrow \D A$ by taking
$p_A u = u(1)(x)$. Hence we get an element of $R\rightarrow \isEquiv(\eta_A)$, and so of
$\norm{R}\rightarrow\isEquiv(\eta_A)$
which provides a factorization of a coherently constant map $R\rightarrow A$ through
$R\rightarrow\norm{R}$.

\subsection{Case of a monoid} \label{subsec:monoid}

 We consider the special case where the base category is a monoid $M$.
 If $\vec{x}$ is a sequence $(x_0,\dots,x_n)$,
 we write $t_k\vec{x}$ for the sequence where we omit $x_k$ and replace
 $x_{k+1}$ by $x_kx_{k+1}$ for $k<n$
 and $t_n\vec{x}$ is the sequence where we omit $x_n$. A type in the presheaf
 model is a type $A$ with an $M$\hyp{}action, and an element of $\D A$ is then
 a family of elements $u(\vec{i})(\vec{x})$ in $A$ with $\vec{i}$ in $\PP_n$
 and $\vec{x}$ in $M^{n+1}$ satisfying the compatibility conditions
 \begin{enumerate}
   \item $u(\vec{i})(\vec{x}) = u(s_k\vec{i})(t_k\vec{x})$ on $i_k = 0$ for $k<n$
     and
   \item $u(\vec{i})(\vec{x}) = u(s_n\vec{i})(t_n\vec{x})x_n$ on $i_n = 0$
     \end{enumerate}
 We define the $M$\hyp{}action on $\D A$ by
 $ux(\vec{i})(x_0,\dots,x_n) = u(\vec{i})(xx_0,\dots,x_n)$.

As a special case, let $M$ be the walking idempotent.
Let $e^2 = e$ be the non-trivial idempotent element of $M$.
Here is an example of a non\hyp{}modal type which is pointwise contractible, but not contractible.
Let $\Gamma$ be the set with elements $\rho_1,\rho_2$ and $\rho$ with
$\rho_1e = \rho_2 e = \rho$. We let $A$ be the following type.
We let $A\rho_1$ be the point $a_1$ and $A\rho_2$ be the point $a_2$
and $A\rho$ be the groupoid with two isomorphic objects $u_1,u_2$
with $a_ie = u_i$ for $i=1,2$. The type $A$ is then pointwise contractible
but it has no global point.%
\footnote{If $a$ is such a point, we should
have $a\rho_i = a_i$ and then $(a\rho_i) e = u_i$ and
$a(\rho_1e) = a(\rho_2e) = a\rho$ which is not possible since $u_1,u_2$
are distinct.}

\subsection{Generalization to a Grothendieck topology} \label{subsec:gen-grothendieck-top}

A Grothendieck topology $\JJ$ on the category $\CC$ defines
a set $\Cov(X,I) = \JJ(X)$ and we have a family
$\E_S$ indexed by $S:\Cov$ defined as follows. Let $\rho$ be in $\Gamma(X)$,
and $S$ is in $\Gamma\rightarrow\Cov$, so that
$S\rho$ is in $\Cov(X) = \JJ(X)$, which is a set of sieves on $X$.

 An element
of $(\E_SA)\rho$ is now a family $u(f)$ in $A\rho f$
{\em with $f$ in $S\rho$}.
We define in this way a family of lex operations $\E_S$
and an associated family of descent data operations $\D_S$
indexed by $S:\Cov$.

Note that if $S_1\rho$ is a subset of $S_2\rho$ for all $\rho$,
then we have a canonical projection map $\D_{S_2}A\rightarrow \D_{S_1}A$ that coheres with the pointings.
If $A$ is $\D_{S_1}$\hyp{}modal a left inverse of $\eta^{S_1}_A$ composed with this projection
map is a left inverse of $\eta^{S_2}_A$. Hence a  $\D_{S_1}$\hyp{}modal type
is also  $\D_{S_2}$\hyp{}modal and we have $\D_{S_1}\leqslant \D_{S_2}$ for the preorder defined
in \cref{Suniv}.
Since $\JJ$ is a Grothendieck topology, the family $\D_S$ over $S : \Cov$ is filtered.
Thus, we can apply \cref{filter} to obtain a model of univalent type theory with higher inductive types.
This can be seen as constructively modelling higher sheaves over $\JJ$ in the cubical model over $\BB$.

The next \lcnamecref{Gcontr} will be used for building such a sheaf model where
countable choice does not hold. The proof is similar to the one of \cref{Dfill}.

\begin{proposition}\label{Gcontr}
  If $A$ in $\Type(\Gamma)$ and $S$ in $\Gamma\rightarrow\Cov$ and $A$ is $D_S$-modal
  and $\rho$ in $\Gamma(X)$
  and $A\rho f$ is (pointwise) contractible for each $f$ in $S\rho$
  then we can find a uniform extension operation $e_{A\rho}(f,\psi,u)$ in $A\rho f$
  for all $f:Y\rightarrow X$ and $u$ partial element in $A\rho f$ of extent $\psi$.
\qed
\end{proposition}

By uniform, we mean that we have
%
\begin{equation*}
e_{A\rho}(f,\psi,u)g = e_{A\rho}(fg,\psi,ug)
\end{equation*}
%
in $A\rho fg$ for any $g:Z\rightarrow Y$.

\subsection{A model with the negation of countable choice}

%% \fxnote{Christian:
%% This subsection was using HoTT language (\eg proposition means homotopy proposition).
%% This feels very appropriate, since it is about statements in HoTT.
%% But it conflicts with our standing assumption that ``strict'' is the default.
%% So I changed it.
%% An alternative would to be make an exception just for this subsection.
%% }

  Using in an essential way the notion of {\em homotopy} descent data, we build a model
  of univalent type theory with higher inductive types
  with a countable family of sets $\E_n$ such that each the homotopy propositional truncation $\norm{\E_n}$ is inhabited,
  but $\norm{\prod_{n:N} \E_n}$ is not globally inhabited.

We consider the following space, corresponding to the lattice generated by formal elements $X_n$
and $L_n$ with the relations $X_0 = 1$, $X_n = L_n\vee X_{n+1}$ and $L_{n+1} = L_n\wedge X_{n+1}$.
%where $X_n$ is covered by $L_n$ and $X_{n+1}$
%We then define $A_0$ and $A_1$ as follows: we take $A_1(n)$ to be $1$ on $L_0$ while
%$A_0(n)$ is $1$ on $X_n$.
%where we only show $A_0,A_1,A_2$
Using \cref{Gcontr} one can show the following result.

\begin{proposition}\label{counter}
  The type $\norm{L_0+X_n}$ is contractible for all $n$
  while $L_0$ is the homotopy propositional truncation of $\prod_{n:N} (L_0 + X_n)$.
%  Each $\norm{A_n}$ is the true proposition, while $B = \prod_{n:N} A_n$ is the strict proposition
%  with $B(X_n) = \emptyset$ and $B(L_n) = \{0\}$.
\qed
\end{proposition}

\begin{corollary}
  There exists a model of univalent type theory with higher inductive types where countable choice does not
  hold.
\qed
\end{corollary}

As stressed in~\cite{SwanU19}, it is yet unknown how to build a model of univalent type
theory and higher inductive types satisfying countable choice in a constructive metatheory.
(Countable choice holds in a classical metatheory in the simplicial set model.)

%% \subsection{Equivalences between general presheaves}

%% We have a fibrant replacement operation $\overline{A}$ with a trivial cofibration
%% $A\rightarrow \overline{A}$. Also at each level $X$, we have a fibrant replacement
%% $A(X)\rightarrow \overline{A(X)}$. There is a canonical map $\overline{A(X)}\rightarrow \overline{A}(X)$
%% and this is an equivalence between two fibrant types at level $X$.

%% A map $\sigma:A\rightarrow B$ determine a map $\overline{\sigma}:\overline{A}\rightarrow \overline{B}$
%% and $\sigma$ is an equivalence if and only if $\overline{\sigma}$ is an equivalence. By \cref{equiv},
%% this is the case if and only if each map $\overline{A}(X)\rightarrow \overline{B}(X)$ is an equivalence at level $X$,
%% which in turn holds if and only if each map
%% $\overline{A(X)}\rightarrow \overline{B(X)}$ is an equivalence at level $X$, and hence if and only if
%% each map $A(X)\rightarrow B(X)$ is an equivalence.

%\section{Type\hyp{}theoretic formulation}

\section{Variation with another notion of cofibration} \label{sec:variation}

We explain how to modify the definition of filling operation if we
work with the other notion of cofibration classified by $\cofunivLW$.
Recall that an element of $\cofunivLW(X)$
is no longer constant, but
is given by a family of elements  $\psi_f$ in $\cofuniv_{\BB}$ for $f:Y\rightarrow X$ and
satisfying $\psi_{f}\leqslant \psi_{f g}$ if $g:Z\rightarrow Y$.

All the main results above still hold for this new notion of cofibration,
suitably modified.
The notion of {\em filling operation}
for $A$ is given by an operation $c_A$ which takes as
argument $\gamma$ in $\Gamma(X)^{\II_{\BB}}$ and $\psi$ in $\cofuniv_\BB(X)$
and a family of elements $u(i)$ in $A\gamma(i)f$ on the extent
$\psi_f \vee i = 0$ such that $u_f(i)g = u_{fg}(i)$  for $g:Z\rightarrow Y$
on the extent
$\psi_f \vee i = 0$.
(There is a dual operation with $i=1$ instead.)
It produces an element $c_A(X,\gamma,\psi,u)(i)$ in $A\gamma(i)$
such that

\begin{enumerate}
  \item $c_A(X,\gamma,\psi,u)(i)f = u_f(i)$ on $\psi_f\vee i=0$,
  \item $c_A(X,\gamma,\psi,u)(i)f = c_A(Y,\gamma',\psi f,u')(i)$ with
    $\gamma'(i) = \gamma(i)f$ and $u'_g(i) = u_{fg}(i)$ on the extent
    $\psi_{fg}\vee i=0$ for $g:Z\rightarrow Y$.
\end{enumerate}

For instance, \cref{Efill} becomes the following result.

\begin{lemma}
  If $A$ has a pointwise filling operation $c_A(X)$ then $\E A$ has a filling operation.
\end{lemma}

\begin{proof}
  We take $\gamma$ in $\Gamma(X)^{\II_{\BB}}$ and $u_f(i)$ in $(\E A)\gamma(i)f$
  on the extent $\psi_f\vee i = 0$ and we define $v(i) = c_{\E A}(X,\gamma,\psi,u)(i)$
  in $(\E A)\gamma(i)$. For $f:Y\rightarrow X$, we take (filling at level $Y$)
  %
  \begin{equation*}
  v(i)(f) = c_A(Y)(\gamma',\psi',u')
  \end{equation*}
  %
  where $\gamma'(i) = \gamma(i)f$ and $\psi' = \psi_f$
  and $u'(i) = u_f(i)(\id_Y)$ in $A\gamma(i)f$ on the extent $\psi_f\vee i=0$.
\end{proof}

Let us give some examples.

\medskip

The first example is when $\CC$ is the poset $0 \leqslant 1$.
In this case, a global type $A$ is given by two spaces with a map $A(1)\rightarrow A(0)$.
An element of $\cofunivLW(0)$ is an element of $\cofuniv_{\BB}$ while an element of $\cofunivLW(1)$ is a pair $\psi_1,\psi_0$ of elements of $\cofuniv_{\BB}$ with $\psi_1\leqslant\psi_0$.
One can check that $A$ is fibrant exactly if $A(0)$ is fibrant and $A(1) \to A(0)$ is a fibration, and a similar characterization holds in the relative situation (for a type $A$ over $\Gamma$) and for trivial fibrations.
Using \cref{all-modal-via-pointwise-contractible:contractible} of \cref{all-modal-via-pointwise-contractible}, one sees that every type in the model is $\D$\hyp{}modal.
The model coincides with the Reedy presheaf model described in~\cite{Shulman15} over the direct category $\CC$ in the model of univalent type theory given by the base model.
More generally, this will be the case for an arbitrary direct category $\CC$ for which the inclusion of objects into morphisms given by identities is decidable.

\medskip

The second example is the walking retract $\CC$ generated by maps $f \colon 0 \to 1$ and $g \colon 1 \to 0$ satisfying $g f = \id_0$.
Note that $\CC$ is the idempotent splitting of the walking idempotent monoid $\M$ considered in \cref{subsec:monoid}.
This makes the cubical presheaf models (for both $\cofunivConst$ and $\cofunivLW$) over $\CC$ and $\M$ equivalent.
Level $0$ in the model over $\CC$ correspond to the fixpoints of the action of $e$ in the model over $\M$.
Taking $\cofunivLW$ as the cofibration classifier, the model of modal types gives a model for pointed families in a cubical model.
It is homotopically correct in the sense that the equivalences are levelwise.

\medskip

One might ask if types in the above model are already $\D$\hyp{}modal, similar to what happens for the poset $0 \leqslant 1$.
More generally, one might attempt to generalize from a direct category $\CC$ to a Reedy category $\CC$ that is elegant~\cite{BergnerR13}; the walking retract is an example of an elegant Reedy category, with coface map $f$ and codegeneracy map $g$.
Taking $\cofunivLW$ as the cofibration classifier, one might ask if the (trivial) fibrations are given by the (trivial) Reedy fibrations; as before, this would imply that every type in the model is $\D$\hyp{}modal.
An equivalent condition is that the levelwise cofibrations (classified by $\cofunivLW$) are also the Reedy cofibrations.
This holds true in classical situations where cofibrations and monomorphisms coincide and gives rise to the classical model~\cite{Shulman13} over an elegant Reedy category.

Unfortunately, this fails to hold in our constructive setting.
Ultimately, this is because the inclusions $A(X) \to A(Y)$ are not generally cofibrations for a global type $A$ and a codegeneracy map $Y \to X$ in $\CC$.
For the case of the walking retract, this is the inclusion $A(0) \to A(1)$.
In terms of a global type $A$ in the model over the walking monoid $\M$, it is the inclusion of fixpoints of the action of $e$ on $A$.
For a counterexample, let $S$ be a discrete space with non-decidable equality in one of the concrete cubical models listed in~\cref{subsec:examples}.
Take $A = S \times S$ with the action of $e$ given by swapping.

\section{Related and future work}

Shulman~\cite{Shulman19} shows that all $(\infty,1)$\hyp{}toposes have strict univalent universes, using a classical metatheory.
This work does not cover however (yet) higher inductive types and cumulativity of universes.
There are close connections between Shulman's work and ours, which we plan to explore in future work.
His work inspired some results about pointwise weak equivalences in~\cref{sec:homotopy-descent-data}, in particular~\cref{Dequiv}.

Once we have a presheaf model of univalence with homotopical features such as ours, it is now understood (see \eg~\cite{Sattler17,Boulier18}) how to define a Quillen model structure whose (trivial) fibrations coincide with the (contractible) types.
For the model of $\D$\hyp{}modal types, we expect that, similar to~\cite{Shulman19}, that the weak equivalences are the levelwise weak equivalences and the fibrations are a variation%
\footnote{We define a family of types to be injectively fibrant if it lifts against cofibrations that are levelwise trivial cofibrations.}
of the injective fibrations.
We leave this to future work.

%Recall the parameter $\CC$ to our theory, an external category.
Instead of parameterizing our construction over an external category $\CC$, we could start from a internal category $\CC$ in presheaves over $\BB$.
Note that the category of presheaves over an internal category in presheaves is still a presheaf category.
Compared to the construction of~\cite{Shulman19} (which instantiates at this level of generality), we seem to need less fibrancy assumptions on this internal category.
We leave this generalization to future work.

%Connections with \cite{Anel20} should be explored.
%\fxnote{This seems really far off compared to Shulman's work.}

% Christian: seems a bit far fetched
%The ultimate goal, following the key insights of~\cite{Voevodsky15,HoTT}, is to rewrite
%and develop higher topos theory in a constructive and type\hyp{}theoretic way.

\section*{Acknowledgements}

Many thanks to Mathieu Anel, Steve Awodey, Mart\'{\i}n Escard\'{o}, Eric Finster,
Dan Licata, Emily Riehl, Mike Shulman, Bas
Spitters and Matthew Weaver for many discussions and remarks.

\appendix

\begin{comment}
\section{Categorical analysis of the construction} \label{sec:injective-fibrations}

The purpose of this \lcnamecref{sec:injective-fibrations} is twofold.
First, we give an alternative, categorical presentation of the homotopy descent data operation constructed in \cref{sec:homotopy-descent-data}.
We aim to be abstract, working at a level of generality that includes the variant of \cref{subsec:gen-grothendieck-top} and other cases.
We give an alternative proof that the homotopy descent data operation indeed forms a modality that sidesteps the combinatorics of\cref{key2}.

Second, we use this to connect our construction to recent work by Michael Shulman~\cite{Shulman19}.
In particular, we will see that our homotopy descent data operation is a semisimplicial version of the cobar construction in~\cite{Shulman19} in a non\hyp{}simplicial setting.
In contrast to~\cite{Shulman19}, the instance of the cobar construction here forms a lex modality in an ambient larger model.
This greatly simplifies the construction of the model for higher presheaves, which is then simply given by localization at this lex modality.

Following~\cite{Shulman19}, we define an analogue of injective fibrations in our setting.
Under classical assumptions, these coincide with injective fibrations as usually considered.
We show that the modal types of the homotopy descent data operation are in correspondence with injective fibration.

\subsection{Homotopical preliminaries in a cubical model} \label{subsec:homotopical-preliminaries}

Before, we can come to the main purpose of this \lcnamecref{sec:injective-fibrations}, we must recall some basic homotopical notions in a cubical model.
Given a small category $\mathcal{A}$, we write $\widehat{\mathcal{A}}$ for the category of presheaves over $\mathcal{M}$.
We now fix a cubical model $\widehat{\BB}$ as in \cref{cubical-sets}.
Let us review basic categorical facts about this cubical model.
We will later use these also for cubical models of presheaves in $\widehat{\BB}$.

We follow~\cite{GambinoSattler17} in setting up basic algebraic homotopical notions.
We start by defining relevant categories of arrows in $\widehat{\BB}$ (\ie, categories over $\widehat{\BB}^\to$).
This makes use of the lifting operation for categories of arrows~\cite[Paragraph~3.8]{Garner09}.
Two categories $\mathcal{S}$ and $\mathcal{T}$ of arrows are called \emph{logically equivalent}, denoted $\mathcal{S} \leftrightarrow \mathcal{T}$, if we have maps in both directions between them.
In many situations, we can treat logically equivalent categories of arrows as interchangeable.
\begin{itemize}
\item
The category of \emph{generating cofibrations} $\GenCof$ is the (locally small) category of arrows classified by the cofibration classifier $1 \to \cofuniv$, with pullback squares as morphisms.
It has a small replacement $\GenCof' \to \GenCof$, the restriction to arrows targeting representables.
Here, we use the word ``replacement'' in the sense that the canonical map $(\GenCof')^\pitchfork \to \GenCof^\pitchfork$ is iso.
\item
The category of \emph{generating trivial cofibrations} $\GenTrivCof$ is the ``box pushout product'' of $\GenCof$ with the interval endpoint inclusions $0, 1 \colon 1 \to \II$.%
\footnote{%
Pushout products with cofibrations such as $0, 1 \colon 1 \to \II$ exist in $\widehat{\BB}$ without assuming quotients in the metatheory.
Instead one uses that cofibrations are levelwise decidable inclusions to implement these pushouts in $\widehat{\BB}$ using coproducts at every level.
}
Note that its maps still form pullback squares.
If we are in the setting of~\cite{AngiuliBCHHL} instead of~\cite{OrtonP16}, we instead define $\GenTrivCof$ as the ``box pushout product'' in $\widehat{\BB}_{/\II}$ of $\GenCof_{/\II}$ with the generic point $\II^*(1) \to \II^*(\II)$ followed by the forgetful functor $\widehat{\BB}_{/\II} \to \widehat{\BB}$.
We obtain a small replacement $\GenTrivCof' \to \GenTrivCof$ by starting with $\GenCof'$ instead of $\GenCof$.
\item
The category of \emph{trivial fibrations} is $\TrivFib = \GenCof^\pitchfork$.
We may describe trivial fibrations with target $X$ using the $\cofuniv$\hyp{}partial map classifier $P_X$~\cite[A2.4]{Johnstone02a} over $X$.
%This is the pointed endofunctor given by pushforward along $X^*(\top) \colon 1 \simeq X^*(1) \to X^*(\cofuniv)$ followed by the forgetful functor.
In the internal language of $\widehat{\BB}$, for a type $Y$ in context $X$, this is $(P_X Y)(x) = \textstyle\sum_{\psi : \cofuniv} Y(x)^\psi$ with pointing sending $y$ to $(\top, \const(y))$.
We call
\begin{equation} \label{partial-map-factorization}
\begin{gathered}
\xymatrix@C+0.5cm{
  Y
  \ar[r]^-{\CofFun'(f)}
&
  P_X Y
  \ar[r]^-{\TrivFibFun'(f)}
&
  X
}
\end{gathered}
\end{equation}

the \emph{$\Phi$\hyp{}partial map factorization} of $f \colon Y \to X$.
This is a stable functorial factorization $(\CofFun', \TrivFibFun')$.%
\footnote{A functorial factorization is \emph{stable} if it is preserved under base change.
In detail, its functorial action on a morphism of arrows that is pullback square consists of two pullback squares.
}
A trivial fibration structure on $Y \to X$ corresponds to a section of $Y \to P_X Y$ over $X$.
We have called this an \emph{extension operation} in \cref{subsec:internal-language-description}.
It allows us to extend any partial element of $Y(x)$ to a total element.
It is logically equivalent to the structure of a contractible type.
More generally, $\TrivFib$ is isomorphic to the category of algebras of the pointed endofunctor $\TrivFibFun'$.
\item
The category of \emph{fibrations} is $\Fib = \GenTrivCof^\pitchfork$.
Using adjointness, the structure of a fibration on a map $Y \to X$ corresponds to sections of the pullback exponentials of $Y \to X$ with the interval endpoints.
We have called this a \emph{filling operation} for $Y$ in context $X$ in \cref{subsec:internal-language-description}.
It is logically equivalent to the structure of a composition operation, which is usually taken to define the types in a cubical model.
If we are in the setting of~\cite{AngiuliBCHHL} instead of~\cite{OrtonP16}, it instead corresponds by adjointness to a section in the slice over $\II$ of the pullback exponential of $Y \to X$ with the generic point $\II^*(1) \to \II^*(\II)$.
This is called \emph{Kan filling} for $Y$ in context $X$ in~\cite{AngiuliBCHHL}, there equivalent to \emph{Kan composition}, and again defines the types of the model.
\item
The category of \emph{cofibrations} is $\Cof = {}^\pitchfork \TrivFib$.
We note the following about the $\Phi$\hyp{}partial map factorization~\labelcref{partial-map-factorization}:
\begin{itemize}
\item
$\CofFun'(f)$ is naturally a pullback of the universal generating cofibration $\top \colon 1 \to \Phi$.
This shows the left functor of the factorization is functorially valued in generating cofibrations $I$.
\item
Using vertical composition of generating cofibrations, \ie the assumption that the maps classified by $\cofuniv$ form a dominance~\cite{Rosolini86}, the pointed endofunctor $\TrivFibFun'$ extends to a monad (lying over the identity monad with respect to the codomain functor).
This shows the right functor of the factorization is functorially valued in trivial fibrations $\TrivFib$.
\end{itemize}
By the functorial retract argument, it follows that $I \leftrightarrow \Cof$, \ie the categories of generating cofibrations and cofibrations are logically equivalent.
In many situations, this justifies not making a distinction between cofibrations and generating cofibrations.
Note that all objects are (functorially) cofibrant.
\item
The category of \emph{trivial cofibrations} is $\TrivCof = {}^\pitchfork \Fib$.
\end{itemize}

Recall the notion of algebraic weak factorization system (awfs) developed in~\cite{Garner09} (there called natural weak factorization system following~\cite{RosickyT07}).
We denote an awfs by its factorization functor $(L, R)$.
An awfs $(L, R)$ has \emph{left} and \emph{right categories} of maps $\Map(L)$ and $\Map(R)$, given by the copointed endofunctor coalgebras of $L$ and pointed endofunctor algebras of $R$, respectively.
Their underlying classes of maps form an ordinary weak factorization system (wfs).
Sometimes, we will have use for the categories $\Coalg(L)$ and $\Alg(R)$ of (co\discretionary{-)}{}{)}monad algebras.

We use the algebraic small object argument~\cite{Garner09} to generate awfs's for our notions of (trivial) fibration.
Note that following~\cite{Swan18b} these awfs's may be defined constructively and without quotients.
Ultimately, this is due to the fact that the universal cofibration $\top : 1 \to \cofuniv$ is a levelwise decidable inclusion.
\begin{itemize}
\item
We have an awfs $(\CofFun, \TrivFibFun)$ algebraically free on $\GenCof$ and $\GenCof'$.
Its left category is $\Cof$ and its right category is logically equivalent to $\TrivFib$.
Separately, we also have the following (see~\cite[Remark~9.5]{GambinoSattler17}).
By~\cite[Section~4]{BourkeG16a}, the $\Phi$\hyp{}partial map factorization $(\CofFun', \TrivFibFun')$ extends to an awfs such that $\Coalg(\CofFun')$ and $\GenCof$ are isomorphic double categories of squares.
Its left category is logically equivalent to $\Cof$ (because $\Cof$ is functorially closed under retract) and its right category is $\TrivFib$.
\item
We have an awfs $(\TrivCofFun, \FibFun)$ algebraically free on $\GenTrivCof$ and $\GenTrivCof'$.
Its left category is $\TrivCof$ and its right category is logically equivalent to $\Fib$.
\end{itemize}
Note that we have a map $\GenTrivCof \to \GenCof$ (in the case of~\cite{AngiuliBCHHL}, this uses ``diagonal  cofibrations'', \ie the assumption that $\II \to \II \times \II$ is a cofibration).
This induces a morphism
\[
(\TrivCofFun, \FibFun) \longrightarrow (\CofFun, \TrivFibFun)
\]
of awfs's.
In particular, we obtain maps $\TrivCof \to \Cof$ and $\TrivFib \to \Fib$.

\begin{remark} \label{cylinder-enriched}
The two awfs's are compatible with the interval in the usual (here algebraized) homotopical sense.
We write $[0, 1] \colon \partial \II \to \II$ for the interval boundary inclusion.
\begin{itemize}
\item
Pushout product with the interval endpoint inclusions $0, 1 \colon 1 \to \II$ lifts to a functor $\Cof \to \TrivCof$.
\item
Pushout product with $\partial \II \to \II$ lifts to a functor $\Cof \to \Cof$.
\item
Pushout product with $\partial \II \to \II$ lifts to a functor $\TrivCof \to \TrivCof$.
\end{itemize}
All of these properties have equivalent dual formulations on the right, \ie using (trivial) fibrations.

Product with the interval defines a notion of (elementary) \emph{homotopy}.
From this, we obtain a notion of \emph{homotopy equivalence}.

More generally, the pushout product with a cofibration preserves (trivial) cofibrations (functoriality on both sides).
As opposed to the above cylinder functor properties, the case of trivial cofibrations does not descend to slices.
\end{remark}

Fibrant univalent universes and dependent products in the cubical model $\widehat{\BB}$ give rise to the fibration extension property and the Frobenius property for the wfs (trivial cofibration, fibration).
Applying~\cite[Theorem~2.6]{Sattler17}, these may be used to show that the wfs's (cofibration, trivial fibration) and (trivial cofibration, fibration) form a model structure.%
\footnote{
There is a subtlety if one does not wish to assume that there are arbitrarily large Grothendieck universes.
In that case, one has to show extension of fibrations along trivial cofibrations by hand, without taking the shortcut of fibrancy of a universe.
For this, one reduces to the case of a generating trivial cofibration by inducting over ``algebraic cell complexes''.
This is a descent argument, using that all colimits involved in the transfinite (in concrete situations, usually countably infinite) construction of \cite{Kelly80} are van Kampen (note that this is not the case for the ``explicit form'' provided in Section~17 of \ibid).
The case of a generator follows from the equivalence extension property, which can be proven without reference to universes.
}
Thus, we obtain an \emph{algebraic model structure} (ams) in the sense of Riehl~\cite{Riehl11}.
This ams is furthermore monoidal (in the algebraic sense).

We will be interested in lifting homotopical structures such as the above wfs's and model structure in $\widehat{\BB}$ to the functor category $[\mathcal{C}^\op, \widehat{\BB}]$ for any category $\mathcal{C}$.
As observed by Riehl~\cite{Riehl11}, this is always possible for awfs's in a functorial manner.
This selling point of awfs's is due to the fact that they are a completely algebraic notion (as opposed to wfs's, which are defined via lifting properties).
However, Emily Riehl's notion of algebraic model structure is insufficiently algebraic to do this job in the world of model structures.
For this reason, we instead use Andrew Swan's notion of ams with structured weak equivalences~\cite[Definition~4.3]{Swan18c}, abbreviated to \emph{structured ams}.
It is fully algebraic and hence lifts to functor categories in a functorial manner.

\begin{remark} \label{algebraic-model-structure}
We can go beyond the model structure and show that our awfs's form a structured ams.%
\footnote{This is perhaps remarkable because the Kan model structure on simplicial sets does not seem to extend to a structured ams.}
We denote the category of weak equivalences by $\WEquiv$.
This can be obtained from an algebraic version of the approach of~\cite{Sattler17}.
For this, one strengthens the fibration extension property (and hence equivalence extension property) and Frobenius property to algebraic operations that are suitably functorial.
Note that fibrancy of the universe and dependent products in the model do give algebraic operations, but they are functorial not with respect to arbitrary morphisms, but only Cartesian morphisms of fibrations (this corresponds to stability under base change).
The functorial Frobenius property is handled in~\cite{GambinoSattler17} (in case of the setting~\cite{OrtonP16}).
The functorial fibration extension property ultimately reduces to checking that the equivalence extension operation satisfies the expected extra functoriality.
We note that Andrew Swan\footnote{Private communication} has suggested an approach that avoids checking all this extra functoriality by hand and instead folds it into a replay of the non\hyp{}algebraic argument within a suitable Grothendieck fibration.
\end{remark}

\begin{remark} \label{model-structure-not-used}
Large parts of our development will not depend on the model structure on $\widehat{\BB}$, or even the wfs (trivial cofibration, fibration).
Arguments and discussions that do will be explicitly marked unless the dependency is obvious.
This is to keep our work accessible to the reader not familiar with the details of the model structure (or even structured ams), in particular the small object argument.
\end{remark}

\begin{remark} \label{slice-cubical-model}
Given a cubical model $\widehat{\BB}$ and $X \in \BB$, then the slice $\widehat{\BB}/X$ again forms a cubical model with
\begin{align*}
\II_{\widehat{\BB}/X} &= X^*(\II_{\widehat{\BB}})
,\\
\cofuniv_{\widehat{\BB}/X} &= X^*(\cofuniv_{\widehat{\BB}/X})
.\end{align*}
Note that $\widehat{\BB}/X$ is equivalent to presheaves over the category of elements of $X$.
Because the partial $\Phi_\BB$\hyp{}map factorization~\labelcref{partial-map-factorization} is stable, the (generating) cofibrations and trivial fibrations in the cubical model $\widehat{\BB}/X$ coincide with those in the slice over $X$ of the cubical model $\widehat{\BB}$.

However, this is not the case for trivial cofibrations and fibrations.
We have a map
\[
\GenTrivCof_{\widehat{\BB}/X} \to (\GenTrivCof_{\widehat{\BB}})_{/X}
,\]
but not generally in the other direction (one case where this happens is when $X$ is $\II$\hyp{}discrete, \ie right orthogonal to $\II \to 1$).
By adjointness, this induces maps
\begin{align*}
(\Fib_{\widehat{\BB}})_{/X} &\longrightarrow \Fib_{\widehat{\BB}/X}
\\
\TrivCof_{\widehat{\BB}/X} &\longrightarrow (\TrivCof_{\widehat{\BB}})_{/X}
.\end{align*}
\end{remark}

\begin{remark} \label{groupoid-nerve}
The category of groupoids $\Gpd$ admits a structured ams with the following categories of maps (up to logical equivalence):
\begin{itemize}
\item
Weak equivalences are equivalences in the structured sense, for example characterized as fully faithful functors (property) with an essentially split epic action on objects (data).
\item
Cofibrations are functors that are decidable inclusions on objects.
This entails that morphisms of cofibrations are pullback squares.
\item
Fibrations are cloven Grothendieck fibrations.
\end{itemize}

When instantiating $\widehat{\BB}$ with one of the standard cubical set models, there is a \emph{nerve} $N \colon \Gpd \to \BB$ induced by the monoidal functor from the chosen cube category to groupoids that sends the ``generator'' to $\braces{0 \simeq 1}$ (if $\widehat{\BB}$ does not contain reversals, one can let the nerve instead start in categories rather than groupoids, with generator sent to $[1] = \braces{0 \to 1}$).
It has a left adjoint, the \emph{fundamental groupoid} functor $\Pi_1 \colon \BB \to \Gpd$.
The adjunction $N \dashv \Pi_1$ is algebraically Quillen.
For this, one easily checks that $\Pi_1$ is an algebraic left Quillen functor.
We will use the nerve $N$ to generate examples of fibrant cubical sets (as \eg~\cref{subsec:examples}).

Since groupoids are functorially fibrant, the nerve functorially preserves weak equivalences (by the functorial version of Ken Brown's lemma).
Note that if $\Phi_\BB$ is the classifier for levelwise decidable inclusions, then the nerve functorially preserves cofibrations, hence also trivial cofibrations.
\end{remark}

\subsection{Homotopical structure in cubical presheaf models}

For a small category $\CC$, we write $\Psh_{\widehat{\BB}}(\CC)$ for the category of presheaves over $\CC$ valued in $\widehat{\BB}$.
This is equivalently the category of presheaves over $\CC \times \BB$ from the previous \lcnamecrefs{sec:presheaf-models}.
Note that $\Psh_{\widehat{\BB}}(\CC)$ is enriched over $\widehat{\BB}$ with the Cartesian monoidal structure.
It is $\widehat{\BB}$\hyp{}bicomplete.
We denote the tensor and power with $K \in \widehat{\BB}$ as $K \times -$ and $(-)^K$, respectively.
These are computed levelwise.

There are two approaches to cubical homotopy theory in $\Psh_{\widehat{\BB}}(\CC)$:
\begin{itemize}
\item
We can lift the structured ams on $\widehat{\BB}$ to presheaves in a functorial manner.
For example, the category of arrows $\Cof$ is given by $\Psh_{\Cof_{\BB}}(\CC) \to \Psh_{\widehat{\BB}^\to}(\CC) \simeq \Psh_{\widehat{\BB}}(\CC)^\to$.
Concretely, the structure of a (trivial) (co\discretionary{-)}{}{)}fibration or structured equivalence on a map $A \to B$ in $\Psh_{\widehat{\BB}}(\CC)$ consists of that structure on $A(X) \to B(X)$ in $\widehat{\BB}^\to$ for every $X \in \CC$ such that the naturality squares of $A \to B$ form morphisms of that structure.
(For this, we note that the relevant categories of arrows $\mathcal{S} \to \widehat{\BB}^\to$ are or can be arranged to be faithful.)
The resulting structured ams is (algebraically) enriched over the structured ams on $\widehat{\BB}$.
\item
We can treat $\Psh_{\widehat{\BB}}(\CC)$ itself as a setting in which to develop a cubical model of type theory as in \cref{subsec:presheaves-in-cubical-sets} and work in the homotopy theory developed in \cref{subsec:homotopical-preliminaries} applied to that setting.
This allows us greater fine\hyp{}tuning of parameters.
Although we will always choose $\II$ as the constant presheaf on $\II_\BB$, we have the two obvious choices $\cofunivConst$ and $\cofunivLW$ for the cofibration classifier.
\end{itemize}
These approaches are related.
\begin{itemize}
\item
The first approach coincides with the choice of the constant cofibration classifier $\cofunivConst$ in the second approach.
This is because all operations involved in defining trivial fibrations (in terms of the partial $\cofunivConst$\hyp{}map classifier) and fibrations (in terms of trivial fibrations) are computed levelwise.
We thus call this the \emph{functorial cubical presheaf model} $\Psh_{\widehat{\BB}}(\CC)_\functorial$.
\item
In the functorial model, the cofibrations are the natural transformations that are levelwise cofibrations with pullbacks as naturality squares.
In contrast, in the model with $\cofunivLW$ as cofibration classifier, the cofibrations are determined levelwise.
More precicely, the category of generating cofibrations is the category of natural transformation that are levelwise generating cofibrations, with pullback squares as morphisms.
We thus call this the \emph{levelwise cofibration cubical presheaf model} $\Psh_{\widehat{\BB}}(\CC)_\levelwise$.
\end{itemize}

\begin{remark} \label{functorial-vs-levelwise}
The levelwise cofibration model has more (trivial) cofibrations than the functorial model.
Thus, it has less (trivial) fibrations.
This is witnessed by maps between categories of maps as appropriate.
In terms of structured ams, the identity lifts to an algebraic left Quillen functor
\[
\Psh_{\widehat{\BB}}(\CC)_\functorial \to \Psh_{\widehat{\BB}}(\CC)_\levelwise
.\]
Let us show that we have a map $\WEquiv_\levelwise \to \WEquiv_\functorial$.
Using functorial fibrant in the levelwise cofibration model and the map $\TrivCof_\levelwise \to \TrivCof_\functorial$, this reduces to the functorial situation of a weak equivalence between fibrant objects in $\Psh_{\widehat{\BB}}(\CC)_\levelwise$.
Any such map is functorially a homotopy equivalence (expressed using the ``cylinder enrichment'' of \cref{cylinder-enriched}) using that the objects are functorially cofibrant\hyp{}fibrant.
But homotopy equivalences are functorially weak equivalences in either of the two models.
\end{remark}

\begin{remark} \label{models-coincide}
If $\CC$ is a groupoid, then the functorial and levelwise cofibration cubical presheaf models coincide.
We thank Emily Riehl for this observation.
This can be seen directly from our characterization of the generating categories of cofibrations.
In particular, the two models coincide if $\CC$ is a discrete category.
\end{remark}

\begin{remark} \label{presheaf-groupoid-nerve}
The nerve\hyp{}realization algebraic Quillen adjunction of \cref{groupoid-nerve} lifts functorially to an algebraic Quillen adjunction
\[
\xymatrix@C+0.5cm{
  \Psh_{\widehat{\BB}}(\CC)_\functorial
  \ar@/^1.2em/[r]^{\textstyle{N}}
  \ar@{{}{ }{}}[r]|-{\textstyle{\bot}}
&
  \Psh_{\Gpd}(\CC)_\functorial
  \ar@/^1.2em/[l]^{\textstyle{\verts{-}}}
\rlap{.}}
\]
Here, the right\hyp{}hand side denotes the functorial lifting of the structured ams from groupoids to groupoid\hyp{}valued presheaves.
It does not seem possible to construct an analogue of the levelwise cofibration model in groupoid\hyp{}valued presheaves.
This is due to the lack of a cofibration classifier in groupoids.
\end{remark}

\begin{remark} \label{weak-equivalences-not-levelwise}
In neither the functorial or levelwise cofibration model, the weak equivalences are characterized levelwise.
For this, we make precise the example from \cref{subsec:examples} based on the groupoid\hyp{}valued presheaf $A$ over $\ints/(2)$ from~\cref{subsec:examples} whose underlying groupoid is cofree on the set $\braces{0, 1}$ with action given by swapping $0$ and $1$.
Let $\widehat{\BB}$ be one of the standard cubical set models with $\Phi_\BB$ classifying all decidable monomorphisms.
By \cref{models-coincide}, the levelwise cofibration and functorial models on $\widehat{\BB}$\hyp{}valued presheaves over $\ints/(2)$ coincide.
By \cref{presheaf-groupoid-nerve}, $N A$ is fibrant and levelwise trivially fibrant in the cubical presheaf model.
But if $N A$ was trivially fibrant, it would have a point.
This would correspond to a point in $A$, which we know does not exist.
\end{remark}

\begin{remark} \label{levelwise-trivial-cofibrations}
In the levelwise cofibration model, one might ask if also the trivial cofibrations are determined levelwise cofibrations).
We do not know the answer to this question.
If it is positive, we could like to call this the left\hyp{}levelwise model.
However, we suspect that the answer cannot be positive in a constructive setting.
\end{remark}

\begin{remark} \label{internal-presheaves}
Instead of a category $\CC$, we can consider a $\widehat{\BB}$\hyp{}enriched category $\CC$ in our presheaf model constructions.
More generally, we can pursue the level of generality of a category $\CC$ internal to $\widehat{\BB}$.
The category of internal presheaves over $\CC$ is the category of presheaves over the externalization of $\CC$.
Thus, it still fits our assumption that cubical models are presheaf categories.
We leave these generalizations to future work.
\end{remark}

\paragraph{Base change}

Let $F \colon \CC' \to \CC$ be a functor of small categories.
$F$ induces a $\widehat{\BB}$\hyp{}enriched restriction functor $F^* \colon \Psh_{\widehat{\BB}}(\CC) \to \Psh_{\widehat{\BB}}(\CC')$.
This functor is $\widehat{\BB}$\hyp{}bicontinuous, in particular preserves tensors and powers.
Kan extension gives adjoints
\begin{equation} \label{base-change-adjunctions}
\begin{gathered}
\xymatrix@C+1.0cm{
  \Psh_{\widehat{\BB}}(\CC')
  \ar@/^2.2em/[r]^{\textstyle{F_!}}
  \ar@/_2.2em/[r]_{\textstyle{F_*}}
  \ar@{{}{ }{}}@/_1.1em/[r]|-{\textstyle{\bot}}
  \ar@{{}{ }{}}@/^1.1em/[r]|-{\textstyle{\bot}}
&
  \Psh_{\widehat{\BB}}(\CC)
  \ar[l]|{\textstyle{F^*}}
}
\end{gathered}
\end{equation}
%
defined by the natural (co\discretionary{-)}{}{)}end formulas
%
\begin{align}
\notag
F_!A(X) &= \int^{X' \in \CC'} \CC(X, F{X'}) \times A(X')  %\colim_{(X', f) \in (X \downarrow F)^\op} A(X')
,\\
\label[formula]{right-kan-extension-formula}
F^*A(X) &= \int_{X' \in \CC'} A(X')^{\CC(F{X'}, X)}  %  \lim_{(X', f) \in (F \downarrow X)^\op} A(X')
\rlap{.}\end{align}

\begin{remark} \label{base-change-exponentiation}
If $F \colon \CC' \to \CC$ is a discrete Grothendieck fibration, then $F$ corresponds via the Grothendieck construction to a (set\hyp{}valued) presheaf over $\CC$.
We can represent this via the constant map $\Set \to \widehat{\BB}$ as an object $\overline{F}$ of $\Psh_{\widehat{\BB}}(\CC)$.
Then $\Psh_{\widehat{\BB}}(\CC')$ is equivalently the slice of $\Psh_{\widehat{\BB}}(\CC)$ over $\overline{F}$.
Under this equivalence, $F^*$ corresponds to pullback $\overline{F}^*$ along the map $\overline{F} \to 1$.
Then $F_!$ corresponds to dependent sum and $F^*$ corresponds to dependent product over $\overline{F}$.
\end{remark}

\begin{remark}[Base change for functorial cubical models] \label{base-change-functorial}
If we endow both sides with the functorial cubical model, then $F^*$ functorially preserves (trivial) (co\discretionary{-)}{}{)}fibrations.
Indeed, $F^*$ forms a morphism of structured ams's.
By adjointness, $F_!$ functorially preserves (trivial) cofibrations and $F_*$ preserves (trivial) fibrations.
Here, we always mean preservation in the algebraic sense, \ie a map of categories of arrows.
\end{remark}

\begin{remark}[Base change for levelwise cofibration models] \label{base-change-levelwise}
Let us also consider the situation where both sides are endowed with the levelwise cofibration cubical model.
We still have $F^* \II \simeq \II$, but $F^*$ preserves the cofibration classifier only laxly in the sense of a pullback square
\begin{equation} \label[square]{levelwise-cofibration-classifier-comparison}
\begin{gathered}
\xymatrix{
  F^* 1
  \ar[r]^{\simeq}
  \ar@{_(->}[d]_{F^* \top}
  \pullback{dr}
&
  1
  \ar@{_(->}[d]_{\top}
\\
  F^* \cofunivLW
  \ar[r]
&
  \cofunivLW
\rlap{.}}
\end{gathered}
\end{equation}
This shows that $F^*$ functorially preserves generating cofibrations.
This also follows directly from our explicit description of the category of generating cofibrations.
Recall that generating trivial cofibrations are constructed from these by operations computed levelwise.
It follows that $F^*$ also preserves generating trivial cofibrations.
By a standard adjointness argument, $F^*$ then functorially preserves (trivial) cofibrations and $F_*$ preserves (trivial) fibrations.

Unfortunately, $F^*$ does not generally preserve (trivial) fibrations.
Let us single out two conditions under which it does (functorially):
\begin{conditions}
\item \label{base-change-levelwise-CC'-groupoid}
$\CC'$ is a groupoid.
Then by \cref{models-coincide} the functorial and levelwise cofibration models on $\Psh_{\widehat{\BB}}(\CC')$ coincide.
By \cref{functorial-vs-levelwise}, (trivial) fibrations in $\Psh_{\widehat{\BB}}(\CC)_\levelwise$ are functorially also (trivial) fibrations in $\Psh_{\widehat{\BB}}(\CC)_\functorial$.
The situation then reduces to \cref{base-change-functorial}.
\item \label{base-change-levelwise-good-exponentiation}
$\CC' \to \CC$ is a discrete Grothendieck fibration.
Then we use \cref{base-change-exponentiation} to describe the adjunction~\labelcref{base-change-adjunctions}.
By \cref{slice-cubical-model}, the equivalence $\Psh_{\widehat{\BB}}(\CC)/\overline{F} \simeq \Psh_{\widehat{\BB}}(\CC')$ functorially sends (trivial) fibrations over $\overline{F}$ in $\Psh_{\widehat{\BB}}(\CC)_\levelwise$ map functorially to (trivial) fibrations in $\Psh_{\widehat{\BB}}(\CC')_\levelwise$.
The situation then reduces to functorial stability of (trivial) fibrations in $\Psh_{\widehat{\BB}}(\CC)_\levelwise$ under pullback.
\end{conditions}
\end{remark}

\mycomment{
When does $F^*$ also preserve (trivial) fibrations (in the algebraic sense)?
Recall that fibrations are characterized by trivial fibrations via operations computed levelwise.
So if $F^*$ preserves trivial fibrations, it also preserves fibrations (both in the algebraic sense).
We are thus reduced to the question of preservation of trivial fibrations.
Inspecting the definition of the partial $\cofunivLW$\hyp{}map classifier, \cref{levelwise-cofibration-classifier-comparison} induces naturally in a map $Y \to X$ in $\Psh_{\BB}(\CC)$ a map $F^* P_X Y \to P_{F^*X} F^* Y$ over $F^* Y$ that makes the square
\[
\xymatrix{
  F^* Y
  \ar[r]^{\id}
  \ar[d]
&
  F^* Y
  \ar[d]
\\
  F^* P_X Y
  \ar[r]
&
  P_{F^* X} (F^* Y)
}
\]
commute.
Given a trivial fibration structure on $Y \to X$, the left map has a retraction.
For a trivial fibration structure on $F^* Y \to F^* X$, the right map should have a retraction.
The question thus reduces to when we can transfer naturally in $Y \to Y$ a retraction of the left map transfers to a retraction of the right map.
For example, this is the case if the map $F^* P_X Y \to P_{F^* X} (F^* Y)$ has a natural retraction.

If $\CC'$ is a discrete category, then the map $F^* P_X Y \to P_{F^* X} (F^* Y)$ has a natural retraction, defined levelwise.
Hence, we get preservation of (trivial) fibrations.
However, it is again simpler to view the problem from a different perspective.
Namely, the functorial and levelwise cofibration models on $\Psh_{\widehat{\BB}}(\CC')$ coincide.
Since (trivial) fibrations in the levelwise cofibration cubical model on $\Psh_{\widehat{\BB}}(\CC)$ are functorially also (trivial) fibrations in the functorial cubical model, preservation of (trivial) fibration reduces to the functorial case.

If $F \colon \CC' \to \CC$ is a discrete Grothendieck fibration, then the map $F^* P_X Y \to P_{F^* X} (F^* Y)$ is invertible.
Hence, we get preservation of (trivial) fibrations.
However, in that case, there is a more conceptual way of seeing this.
Namely, $F$ corresponds to a (set\hyp{}valued) presheaf over $\CC$, which we can represent via the constant map $\Set \to \widehat{\BB}$ as an object $\overline{F}$ of $\Psh_{\widehat{\BB}}(\CC)$.
Then $\Psh_{\widehat{\BB}}(\CC')$ is equivalently the slice of $\Psh_{\widehat{\BB}}(\CC)$ over $\overline{F}$ and $F^*$ is equivalent to the functor of pullback along the map $\overline{F} \to 1$.
This preserves (trivial) fibrations as they are defined by right lifting.
}

\subsection{Abstract perspective on homotopy descent data}

We can now give a more abstract perspective on the lex operation $\E$ and descent data operation $\D$ defined in \cref{sec:a-lex-operation,sec:homotopy-descent-data}.
This will provide connections with~\cite{Shulman19} and allow us to give a different proof that $\D$ forms a modality, sidestepping the key combinatorial ingredient \cref{key2} that established a chain of homotopies between $\eta \D$ and $\D \eta$.

For simplicity of exposition, we only work with functorial cubical presheaf models.
We will remark later on the differences in the situation of levelwise cofibration cubical models.

Let $F \colon \CC' \to \CC$ be a functor of small categories as before.
Again, we consider the base change adjunctions~\labelcref{base-change-adjunctions}.
We focus on $\Psh_{\widehat{\BB}}(\CC)$, defining several new categories of arrows.
Following~\cite[Definition~8.1]{Shulman19}, but working algebraically, we define the categories of \emph{$F^*$\hyp{}cofibrations}, \emph{$F^*$\hyp{}trivial cofibrations}, and \emph{$F^*$\hyp{}weak equivalences} (if we assume structured ams's) in $\Psh_{\widehat{\BB}}(\CC)$ as being created via $F^*$ from the corresponding categories of arrows in $\Psh_{\widehat{\BB}}(\CC')$:
\begin{align*}
\xymatrix{
  F^* \Cof
  \ar[r]
  \ar[d]
  \pullback{dr}
&
  \Cof
  \ar[d]
\\
  \Psh_{\widehat{\BB}}(\CC)
  \ar[r]^{F^*}
&
  \Psh_{\widehat{\BB}}(\CC')
\rlap{,}}
&&
\xymatrix{
  F^* \TrivCof
  \ar[r]
  \ar[d]
  \pullback{dr}
&
  \TrivCof
  \ar[d]
\\
  \Psh_{\widehat{\BB}}(\CC)
  \ar[r]^{F^*}
&
  \Psh_{\widehat{\BB}}(\CC')
\rlap{,}}
&&
\xymatrix{
  F^* \WEquiv
  \ar[r]
  \ar[d]
  \pullback{dr}
&
  \WEquiv
  \ar[d]
\\
  \Psh_{\widehat{\BB}}(\CC)
  \ar[r]^{F^*}
&
  \Psh_{\widehat{\BB}}(\CC')
\rlap{.}}
\end{align*}
We define categories of \emph{injective (trivial) fibrations for $F$} as
\begin{align*}
\InjFib_F &= (\Cof \times F^* \TrivCof)^\pitchfork
,\\
\InjTrivFib_F &= (\Cof \times F^* \Cof)^\pitchfork
.\end{align*}
In particular, an injective fibration is defined as a map that lifts coherently against cofibrations $A \to B$ that become trivial cofibrations under $F^*$ (read algebraically).
Recall that $F^*$ preserves cofibrations.
Thus, $\InjTrivFib_F$ and $\TrivFib$ are logically equivalent.
We have introduced the former only for notational symmetry.

\begin{remark}
The restriction to cofibrations before taking the right lifting closure is necessitated by our technical development.
Our terminology of injective fibration deviates in this aspect from its standard usage.
Note that if we work with levelwise cofibration cubical presheaf models and $F$ is the inclusion $\obj \CC \to \CC$, then $\Cof$ is logically equivalent to $F^* \Cof$.
Then $\InjFib_F \leftrightarrow (\text{$F^*$\hyp{}$\TrivCof$})^\pitchfork$, and the terminology mirrors the classical case.
\end{remark}

\begin{remark} \label{injective-fibrant-is-fibrant}
Since $F^*$ functorially preserves trivial cofibrations, $F$\hyp{}injective fibrancy is stronger than ordinary fibrancy, independently of the choice of $F$.
That is, we have maps $\InjFib_F \to \Fib$ and $\InjTrivFib_F \to \TrivFib$.
These become logical equivalences if $F$ is an equivalence.
\end{remark}

\begin{remark} \label{injective-fibration-pullback-exponential-with-cofibration}
The pullback exponential with a cofibration preserves $F$\hyp{}injective fibrations and $F$\hyp{}trivial injective fibrations, functorially on both sides.
By adjointness, this reduces to functorial closure of $\Cof$, $F^* \TrivCof$, $F^* \Cof$ under pushout product with cofibrations.
Since $F^*$ functorially preserves cofibrations, this reduces to functorial closure of (trivial) cofibrations under pushout product with cofibrations as in \cref{cylinder-enriched}.
\end{remark}

Note that $F_*$ functorially sends (trivial) fibrations to injective (trivial) fibrations for $F$.
This reduces by adjointness to checking that $F^*$ functorially preserves (trivial) cofibrations (\cref{base-change-functorial}).
We now consider the monad $(\E_F, \eta, \mu)$ on $\Psh_{\widehat{\BB}}(\CC)$ induced by the adjunction $F^* \dashv F_*$; explicitly, we have
\[
\E_F = F_* F^*
.\]
It follows that $\E_F$ functorially sends $F^*$\hyp{}fibrations (respectively, $F^*$\hyp{}trivial fibrations) to injective (trivial) fibrations for $F$.
Because both $F_*$ and $F^*$ are $\widehat{\BB}$\hyp{}enriched right adjoints, $\E_F$ preserves $\widehat{\BB}$\hyp{}limits, in particular pullbacks and powers.

\begin{remark}[Lex operation induced by $\E_F$] \label{E-induced-lex-operation}
As a pullback\hyp{}preserving and functorially fibration\hyp{}preserving pointed endofunctor, $\E_F$ internalizes to a lex operation in the cubical model $\Psh_{\widehat{\BB}}(\CC)$.
\fxnote{Add reference to semantic criterion for lex operation from \cref{subsec:lex-operation}.}
To represent the operations on a universe $\UU$ of the model, we need to assume that the Grothendieck universe $U$ used to define $\UU$ contains $\CC$.
In the internal language of $\widehat{\BB}$, we unfold the action of $\E_F$ on a type $A$ in context $\Gamma$.
Using \cref{right-kan-extension-formula}, we obtain the following pullback:
\[
\xymatrix{
  \textstyle\sum_{\gamma : \Gamma(X)} (a_f : A(\gamma f))_{(X', f) \in F \downarrow X}
  \ar[r]
  \ar[d]
  \pullback{dr}
&
  \int^{X' \in \CC'} (\Gamma.A)(F X')^{\CC(F X', X)}
  \ar[d]
\\
  \Gamma(X)
  \ar[r]
&
  \int^{X' \in \CC'} \Gamma(F X')^{\CC(F X', X)}
\rlap{.}}
\]
Here, the family $(a_f)$ is naturally indexed over $f \colon F X' \to X$ with $X' \in \CC'$: for $g \colon Y' \to X'$ a map in $\CC'$, we have $a_f = a_{f \circ F g}$.
If $\CC' \to \CC$ is the inclusion $\obj \CC \to \CC$, we obtain the lex operation $\E$ defined in \cref{sec:a-lex-operation}.
\end{remark}

\mycomment{
The action $\E_F$ on a type $A \in \Type(\Gamma)$ is defined by the following pullback:
\[
\xymatrix{
  \Gamma. \E_F A
  \ar[r]
  \ar@{->>}[d]
  \pullback{dr}
&
  \E_F (\Gamma. A)
  \ar@{->>}[d]
\\
  \Gamma
  \ar[r]^{\eta_\Gamma}
&
  \E_F \Gamma
\rlap{.}}
\]
The action $\EE_F \in \Type(\Gamma. \E_F A)$ on a type $B \in \Type(\Gamma.A)$ is then defined by the following pullback:
\[
\xymatrix{
  \Gamma. A. \EE_F B
  \ar[r]
  \ar@{->>}[d]
  \pullback{dr}
&
  \E_F (\Gamma. A. B)
  \ar@{->>}[d]
\\
  \Gamma. \E_F A
  \ar[r]^{\eta_\Gamma}
&
  \E_F (\Gamma. A)
\rlap{.}}
\]
To represent $\E_F$ and $\EE_F$ as operations on a universe $\UU$ of the model, we need to assume that the Grothendieck universe $U$ used to define $\UU$ contains $\widehat{\BB}, \widehat{\CC}, \widehat{\CC'}$.
Then $F_*$ preserves maps with $U$\hyp{}small fibers.
Since $F^*$ does so trivially, the same is true for $\E_F$.
Let $\UU$ be the universe in the cubical model $\Psh_{\widehat{\BB}}(\CC)$ built from $U$.
Then since $\E_F$ preserves fibrations, the map $\E_F (\sum_{X:\UU} X) \to \E_F \UU$ is again classified by $\UU$ by a map $\E_F \UU \to \UU$.
Precomposing with the unit for $\E_F$, we obtain an internalized version $\UU \to \UU$ of the operation $\E_F$.
}

We write $\Delta$ for the simplex category and $\Delta_+$ for its direct subcategory of face maps.
We use the subscript $\aug$ if we want to talk about the augmented versions.

%We will be interested in homotopy limits of cosemisimplicial diagrams in $\Psh_{\widehat{\BB}}(\CC)$.

\begin{remark}[Functorial model for homotopy limits of semisimplicial shape] \label{weighted-limit-fibrant}
Let $W$ be Reedy cofibrant in $[\Delta_+, \Psh_{\widehat{\BB}}(\CC)]$.
We consider the $W$\hyp{}weighted limit functor
\[
\textstyle\lim^W \colon [\Delta_+, \Psh_{\widehat{\BB}}(\CC)] \to \Psh_{\widehat{\BB}}(\CC)
\rlap{,}\]
taken with respect to the Cartesian monoidal structure.
Explicitly, we have
\[
{\textstyle\lim^W X} = \int^{[n] \in \Delta_+} X_n^{W_n}
\]
It functorially sends levelwise (trivial) fibrations to (trivial) fibrations.
This is because levelwise (trivial) fibrations coincide with Reedy (trivial) fibrations since $\Delta_+$ only has face maps and the pullback weighted limit with a Reedy cofibration functorially preserves Reedy (trivial) fibrations (this reduces to \cref{cylinder-enriched}).

Consider now the model structure on $\Psh_{\widehat{\BB}}(\CC)$.
Let $W$ be a Reedy cofibrant replacement of $1 \in [\Delta_+, \Psh_{\widehat{\BB}}(\CC)]$ (\ie, $W \to 1$ is a weak equivalence).
Then $\lim^W$ provides a functorial model for homotopy limits of levelwise fibrant cosemisimplicial diagrams in $\Psh_{\widehat{\BB}}(\CC)$.
\end{remark}

\begin{remark} \label{weighted-limit-injective-fibrant}
In the setting of \cref{weighted-limit-fibrant}, the $W$\hyp{}weighted limit functor $\lim^W$ functorially sends levelwise $F$\hyp{}injective fibrations to $F$\hyp{}injective fibrations (and the same holds for $F$\hyp{}trivial injective fibrations).
This is proven analogously, but using \cref{injective-fibration-pullback-exponential-with-cofibration} instead of \cref{cylinder-enriched}.
\end{remark}

We now give a conceptual derivation of the operation $\D$ from \cref{sec:homotopy-descent-data}.
The monad $\E_F$ corresponds to a monoidal functor $\E_F^\bullet \colon \Delta_\aug \to \Endo(\Psh_{\widehat{\BB}}(\CC))$.
In particular, for $A \in \Psh_{\widehat{\BB}}(\CC)$, we obtain a cosimplicial diagram
\[
\xymatrix@C+1.0cm{
  \E_F A
  \ar@<+0.8em>[r]^{\eta_{\E_F A}}
  \ar@<-0.8em>[r]_{\E_F \eta_A}
&
  \E_F^2 A
  \ar[l]|{\mu_A}
  \ar@<+1.8em>[r]^{\eta_{\E_F^2 A}}
  \ar@<+0.0em>[r]|{\E_F \eta_{\E_F A}}
  \ar@<-1.8em>[r]_{\E_F^2 \eta_A}
&
  \ldots
  \ar@<+0.9em>[l]%|{\mu_{\E_F A}}
  \ar@<-0.9em>[l]%|{\E_F \mu_A}
\rlap{.}}
\]
We are interested in its homotopy limit, desiring a functorial model for it.
Note that the homotopy limit of a cosimplicial diagram coincides with the homotopy limit of the restriction of the diagram to coface maps.
So to simplify, we throw away the codegeneracy maps, leaving a cosemisimplicial diagram
\[
\xymatrix@C+1.0cm{
  \E_F A
  \ar@<+0.5em>[r]^{\eta_{\E_F A}}
  \ar@<-0.5em>[r]_{\E_F \eta_A}
&
  \E_F^2 A
  \ar@<+1.0em>[r]^{\eta_{\E_F^2 A}}
  \ar[r]|{\E_F \eta_{\E_F A}}
  \ar@<-1.0em>[r]_{\E_F^2 \eta_A}
&
  \ldots
\rlap{.}}
\]
The corresponding functor $V \colon \Psh_{\widehat{\BB}}(\CC) \to [\Delta_+, \Psh_{\widehat{\BB}}(\CC)]$ functorially sends fibrations to levelwise $F$\hyp{}injective fibrations (and trivial fibrations to levelwise $F$\hyp{}trivial injective fibrations).
The coaugmentation provides a natural transformation $\eta \colon \const \to V$.

Let $W$ be a Reedy cofibrant replacement of $1$ in $[\Delta_+, \Psh_{\widehat{\BB}}(\CC)]$ (we can strengthen this to a condition not referring to the model structure by demanding that $W(X) \to 1$ is a homotopy equivalence for each $X \in \CC$ as per \cref{cylinder-enriched}).
Using $\lim^W$ as a strictly functorial model for homotopy limits of cosemisimplicial shape as per \cref{weighted-limit-fibrant}, we define the endofunctor $\D_F$ on $\Psh_{\widehat{\BB}}(\CC)$ as
\[
\D_F = {\textstyle\lim^W} \circ V
.\]
The maps $W \to 1$ and $\eta \colon \const \to V$ provide a pointing $\eta$ of $\D_F$ given by
\[
\xymatrix{
  A
  \ar[r]^-{\simeq}
&
  \textstyle\lim^1 \const(A)
  \ar[r]
&
  \textstyle\lim^W V(A) = \D_F A
\rlap{.}}
\]
As a composite of $\widehat{\BB}$\hyp{}continuous functors, $\D_F$ is again $\widehat{\BB}$\hyp{}continuous.
In particular, like $\E_F$, it preserves pullbacks and powers.
By \cref{weighted-limit-injective-fibrant}, $\D_F$ functorially fibrations to $F$\hyp{}injective trivial fibrations (and trivial fibrations to $F$\hyp{}trivial injective fibrations).

\begin{remark}[Concrete Reedy cofibrant replacements] \label{concrete-reedy-cofibrant-replacements}
Let $\widehat{\BB}$ be one of the standard cubical set models.
Then we can give particularly simple models $W$ for the Reedy cofibrant replacement of $1$ in $[\Delta_+, \widehat{\BB}]$.
These lift functorially to Reedy cofibrant replacements as in \cref{weighted-limit-fibrant}.
\begin{options}
\item \label{concrete-reedy-cofibrant-replacements:P}
We can set $W = \PP$ with $\PP \colon \Delta_+ \to \widehat{\BB}$ from \cref{sec:homotopy-descent-data}.
Recall that $\PP_n$ is the subobject of $\II^{\obj {[n]}}$ of all faces that exclude the initial vertex.
This is evidently Reedy cofibrant; it may be described as the Reedy cofibrant replacement obtained by adding at level $n$ a new ``center'' point and cubes glueing it to the ``boundary'' $L_n \PP$.
Assume now that $\BB$ has $\vee$\hyp{}connection.
Then we can extend $\PP$ to a functor from $\Delta$.
The map $W \to 1$ is a homotopy equivalence; the homotopy inverse $1 \to W$ selects the terminal vertex.
\item \label{concrete-reedy-cofibrant-replacements:Q}
We have a pointed endofunctor $C$ of \emph{cone formation}~\cite{DohertyKLS20} in $\widehat{\BB}$ sending $A$ to the pushout%
\footnote{Since this is a pushout along a cofibration, it can be implemented without metatheoretical quotients.}
\[
\xymatrix{
  \braces{1} \times A
  \ar[r]
  \ar@{>->}[d]
&
  1
  \ar@{>->}[d]
\\
  \II \times A
  \ar[r]
&
  \pullback{ul}
  C A
}
\]
with pointing given by $\braces{0} \times A \to \II \times A$.
This corresponds to a monoidal functor $C^\bullet \colon \Delta_{+,\aug} \to \Endo(\widehat{\BB})$.
Evaluating at $0 \in \widehat{\BB}$, we obtain Reedy cofibrant $\QQ \in [\Delta_+, \widehat{\BB}]$.
Assume now that $\BB$ has $\vee$\hyp{}connection.
Then $C$ extends to a monad, corresponding to $C^\bullet \colon \Delta_{\aug} \to \Endo(\widehat{\BB})$ and restricting to $\QQ \in [\Delta, \widehat{\BB}]$.
Since $\QQ_n A$ is of the form $C B$ for some $B$, the map $\QQ_n A \to 1$ is a homotopy equivalence, although not functorially in $[n] \in \Delta$.
\item \label{concrete-reedy-cofibrant-replacements:R}
%We can construct a simplicial cofibrant resolution using just the Cartesian monoidal structure.
Let $\widehat{\BB}/(\II \times \II)$ be the monoidal category of endospans on $\II$ (the unit is the trivial endospan on $\II$ and composition is given by pullback over $\II$).
The object $M = (\II \times \II, \id)$ has the structure of a monoid: the unit is the diagonal $\II \to \II \times \II$ and the multiplication is the map $\II \times \II \times \II \to \II \times \II$ forgetting the middle copy of $\II$.
This corresponds to a monoidal functor $M^\bullet \colon \Delta_\aug \to \widehat{\BB}/(\II \times \II)$.
We obtain $\RR \in [\Delta, \widehat{\BB}]$ by composing $M$ with pullback along $(0, 1) \colon 1 \to \II \times \II$.

Concretely, we have $\RR_n \simeq \II^n$.
The generators have the following description.
The inner coface maps duplicate coordinates.
The outer coface maps insert $0$ at the beginning and $1$ at the end, respectively.
The degeneracy maps forget coordinates.
It is also possible to describe $\RR$ using the equivalence $\Delta_\aug \simeq \Delta_{\gen}^\op$ where $\Delta_{\gen}$ refers to the wide subcategory of $\Delta$ of \emph{generic} maps, preserving the initial and terminal object.

Reedy cofibrancy of $\RR$ is best checked at the level of $M$ (using that cofibrations are stable under pullback).
The latching object $L_n M$ is the union of all proper generalized $\II$\hyp{}diagonals targetting $M^{[n]} \simeq \II^{n+2}$.
Thus, we need the assumption that the diagonal $\II \to \II \times \II$ is a cofibration.
The map $\RR \to 1$ is levelwise a homotopy equivalence.
If $\BB$ has $\wedge$\hyp{}connections or $\vee$\hyp{}connections, then $\RR \to 1$ into a homotopy equivalence.
We leave this as an exercise to the reader.
\item \label{concrete-reedy-cofibrant-replacements:Riehl-Shulman}
\fxnote{Discuss version of Riehl\hyp{}Shulman (variant of \cref{concrete-reedy-cofibrant-replacements:R})}
\end{options}
\end{remark}

\begin{remark}[Lex operation induced by $\D_F$] \label{D-induced-lex-operation}
Just like $\E_F$ in \cref{E-induced-lex-operation}, so does $\D_F$ induce a lex operation in the cubical model $\Psh_{\widehat{\BB}}(\CC)$ \fxnote{add reference to semantic criterion}.
In addition to the smallness assumption on $\CC$, we also need to assume that $W$ lives in every universe of $\Psh_{\widehat{\BB}}(\CC)$.
The latter is the case using the concrete models of $W$ provided by \cref{concrete-reedy-cofibrant-replacements}.
If $\CC' \to \CC$ is the inclusion $\obj \CC \to \CC$ and $W$ is~\cref{concrete-reedy-cofibrant-replacements:P} from \cref{concrete-reedy-cofibrant-replacements}, we obtain the lex operation $\D$ defined in \cref{sec:homotopy-descent-data}.
The proofs of \cref{Dfill,Dcontr} can be seen as an unfolding of the fibrancy conssiderations in \cref{weighted-limit-fibrant}.
\end{remark}

We want to consider diagrams over $\Delta_{\aug}$ that have extra degeneracies.
This is a diagram over $\Delta_{\bot}$.

Have wide subcategory inclusion $\Delta_\aug \to \Delta_{\bot}$.

\begin{lemma}
Let $X \colon \Delta_{\bot} \to \M$.
Then the restriction $X|_{\Delta} \colon \Delta_\aug \to \M$, seen as a diagram $1 \join \Delta \to M$, is a homotopy limit.
\end{lemma}

\subsection{\texorpdfstring{$\D_F$}{DF}\hyp{}modal types equal injective fibrations} \label{subsec:injective-fibrations}
\end{comment}

\section{General results for lex modalities}\label{sec:general-results}

Some of our results hold for modalities in the sense of~\cite{RijkeSS17} that are not necessarily presented in a strict manner by a lex operation.
The main example is the case of accessible modalities, which are implemented using higher inductive types that rarely give rise to a lex operation.
The purpose of this \lcnamecref{sec:general-results} is to prove these more general statements.
We work in the homotopy type theory setting of~\cite{RijkeSS17}.
Universes are assumed univalent and closed under dependent sums, dependent products, identity types.
For statements involving accessible modalities, we also assume closure under higher inductive types.

In this \lcnamecref{sec:general-results}, we take terminology with potentially both strict and homotopical meaning to have the homotopical meaning by default.
This is opposed to the rest of the article, where we default to the strict meaning.
For example, equality refers to the identity type, and pullbacks refer to homotopy pullbacks (expressed using the identity type).

We write $\Modality(\UU)$ for the type of modalities on a universe $\UU$.
Recall from~\cite{RijkeSS17} that $M : \Modality(\UU)$ has an underlying subuniverse%
\footnote{%
By a subuniverse of $\UU$, we mean a subobject of $\UU$, \ie a predicate on $\UU$.
This is formally a map $\UU \to \Prop$ where $\Prop$ is the universe of (homotopy) propositions.
It is not to be confused with a subuniverse in the set\hyp{}theoretic sense in a model where universes are built out of sets.
We note that the size of the propositions in $\Prop$ here does not matter for us; one choice is propositions in $\UU$, but one could allow also a larger universe.
}
of $\UU$, the $M$\hyp{}modal types $\UU_M$.
Subuniverses of $\UU$ carry an evident poset structure.
Following~\cite[Subsection~3.2]{RijkeSS17}, we obtain a poset structure also on $\Modality(\UU)$.

\begin{definition} \label{modality-extension}
Let $\UU$ be a universe contained in a universe $\UU'$.
A modality $M'$ on $\UU'$ is an \emph{extension} of a modality $M$ on $\UU$ if every $M$\hyp{}modal type in $\UU$ is $M'$\hyp{}modal in $\UU'$ and for $X : \UU$, the canonical map $M'X \to MX$ is invertible.
\end{definition}

The above conditions mean that a $\UU$\hyp{}small type is $M$\hyp{}modal exactly if it is $M'$\hyp{}modal and $M$\hyp{}connected exactly if it is $M'$\hyp{}connected.
In terms of the stable factorization systems $(\mathcal{L}, \mathcal{R})$ and $(\mathcal{L}', \mathcal{R}')$ corresponding to $M$ and $M'$, this means that $\mathcal{L}$ and $\mathcal{R}$ are the restrictions of $\mathcal{L}'$ and $\mathcal{R}'$ to maps between $\UU$\hyp{}small types.
For this, recall~\cite[Subsection~1.2]{RijkeSS17} that the left and right classes of the stable factorization system corresponding to a modality are the connected and modal maps, which are defined by having connected and modal fibers, respectively.

We write $\Modality(\UU < \UU')$ for the type of pairs $(M, M')$ with $M$ a modality on $\UU$ and $M'$ an extension of $M$ to $\UU'$.
The poset structures on $\Modality(\UU)$ and $\Modality(\UU')$ extend to a poset structure on $\Modality(\UU < \UU')$.

The following statement makes precise that up to (essential) size issues, a modality is lex exactly if the universe of modal types is modal.
In particular, a ``size\hyp{}polymorphic'' modality (acting compatibly on all universes) whose action on maps preserves smallness of fibers is lex exactly if universes of modal types are modal.
This generalizes~\cref{univ} to modalities; the smallness condition on fibers mirrors the dependent action $\DD$ on $\UU$\hyp{}small types we require for a lex operation $\D$.
For $M : \Modality(\UU)$, we denote by $\UU_M$ the subuniverse of $\UU$ of $M$\hyp{}modal types.
\begin{proposition} \label{modality-lex-via-universe}
For $(M, M') : \Modality(\UU < \UU')$:
\begin{parts}
\item \label{modality-lex-via-universe:lex-to-universe}
if $M'$ is lex and preserves maps with $\UU$\hyp{}small fibers, then $\UU_M$ is $M'$\hyp{}modal;
\item \label{modality-lex-via-universe:universe-to-lex}
if $\UU_M$ is $M'$\hyp{}modal, then $M$ is lex and $M'$ preserves maps with $\UU$\hyp{}small fibers.
\end{parts}
\end{proposition}

\begin{proof}
For \cref{modality-lex-via-universe:lex-to-universe}, let $M'$ be lex and preserve maps with $\UU$\hyp{}small fibers.
To show that $\UU_M$ is $M'$\hyp{}modal, it suffices to construct a left inverse to $\eta^{M'}_{\UU_M}$ (\cite[Lemma~1.20]{RijkeSS17}).
By univalence of $\UU_M$, this means to find an extension
\[
\xymatrix{
  \textstyle\sum_{X:\UU_M} X
  \ar@{.>}[r]
  \ar[d]
  \pullback{dr}
&
  Z
  \ar@{.>}[d]^{\txt\scriptsize{$\UU$\hyp{}small and\\ $M$\hyp{}modal fibers}}
\\
  \UU_M
  \ar[r]^{\eta^{M'}_{\UU_M}}
&
  M' \UU_M
\rlap{.}}
\]
We use the naturality square of $\eta^{M'}$ at the left map.
The square is a pullback because $M'$ is lex.
The right map has $\UU$\hyp{}small fibers by assumption and has $M$\hyp{}modal fibers because it is $M'$\hyp{}modal as it goes between $M'$\hyp{}modal types.

For \cref{modality-lex-via-universe:universe-to-lex}, assume that $\UU_M$ is $M'$\hyp{}modal.
Then $\UU_M$ is right orthogonal against $M'$\hyp{}connected types, in particular $M$\hyp{}connected types.
This verifies condition~(xiii) of~\cite[Theorem~3.1]{RijkeSS17}, making $M$ lex.
It remains to show that $M'$ preserves maps with $\UU$\hyp{}small fibers.
Given such a map, we factor it using $M$ as an $M'$\hyp{}connected map followed by a map with fibers in $\UU_M$.
Since $M'$ sends $M'$\hyp{}connected maps to equivalences, it remains to show, given $Y : X \to \UU_M$, that $M'(\sum_X Y) \to M'X$ has $\UU$\hyp{}small fibers.
Since $\UU_M$ is $M'$\hyp{}modal, it is right orthogonal against $X \to M'X$.
Thus, $Y : X \to \UU_M$ extends uniquely to a map $Y' : M'X \to \UU_M$.
Looking at the classified maps, we obtain the following commuting diagram:
\[
\xymatrix{
  \textstyle\sum_X Y
  \ar[r]
  \ar[d]
  \pullback{dr}
&
  \textstyle\sum_{z:M'X} Y'(z)
  \ar[d]
\\
  X
  \ar[r]^{\eta^{M'}_X}
&
  M'X
\rlap{.}}
\]
Since the right map has $M$\hyp{}modal (hence also $M'$\hyp{}modal) fibers, it is $M'$\hyp{}modal.
The top map is a pullback of $\eta^{M'}_X$, hence $M'$\hyp{}connected.
Since
\[
\textstyle\sum_X Y \longrightarrow \textstyle\sum_{z:M'X} Y'(z) \longrightarrow M'X
\]
and
\[
\textstyle\sum_X Y \longrightarrow M'(\textstyle\sum_X Y) \longrightarrow M'X
\]
are ($M'$\hyp{}connected, $M'$\hyp{}modal)\hyp{}factorizations of the same map, they coincide.
This shows that the map $M'(\sum_X Y) \to M'X$ is equal to $\sum_{z:M'X} Y'(z) \to M'X$, hence has $\UU$\hyp{}small fibers.
\end{proof}

Recall from~\cite[Subsection~2.3]{RijkeSS17} that accessible modalities admit canonical extensions to larger universes.
If the accessible modality is lex, we observe that it satisfies the technical condition on smallness of fibers of \cref{modality-lex-via-universe}.
This means that \cref{modality-lex-via-universe:lex-to-universe} of that statement can also be regarded as a generalization of the direction from condition~(i) to condition~(iii) in~\cite[Theorem~3.11]{RijkeSS17}.

\begin{corollary}
Let $M$ be an accessible lex modality on a universe $\UU$.
Let $M'$ be its extension to a universe $\UU'$ containing $\UU$.
The $M'$ preserves maps with $\UU$\hyp{}small fibers.
\end{corollary}

\begin{proof}
This follows from \cref{modality-lex-via-universe:universe-to-lex} of \cref{modality-lex-via-universe} since $\UU_M$ is $M'$\hyp{}modal by~\cite[Theorem~3.11]{RijkeSS17}.
\end{proof}

Let $M : I \to \Modality(\UU)$ be a family of modalities.
We write
\begin{equation} \label{subuniverse-structural-meet}
\UU(M) = \textstyle\sum_{X : \UU} \textstyle\prod_{i : I} \,\text{$X$ is $M_i$\hyp{}modal}.
\end{equation}
for the meet of the subuniverses of modal types of $M_i$ over $i : I$.
We call a given meet $\bigwedge M$ of $M$ \emph{structural} if it is preserved under the forgetful functor to the poset of subuniverses.
This means that its subuniverse of modal types is $\UU(M)$.
By~\cite[Theorem 3.11, part (i)]{RijkeSS17}, $M$ has a structural meet exactly if $\UU_M$ admits a reflection in $\UU$.
In that case, $\bigwedge M$ is given by the reflection operation.

Given a family $(M, M') : I \to \Modality(\UU < \UU')$, we say that a given meet of $(M, M')$ is structural if it is sent to structural meets of $M$ and $M'$ by the forgetful functors.
Note that $(M, M')$ has a structural meet exactly if $M$ and $M'$ have structural meets $\bigwedge M$ and $\bigwedge M'$, respectively, and $\bigwedge M'$ is an extension of $\bigwedge M$ to $\UU'$.
This unfolds to the following conditions:
\begin{itemize}
\item
the subuniverse $\UU_M$ of $\UU$ admits a reflection $L$,
\item
the subuniverse $\UU'_{M'}$ of $\UU'$ admits a reflection $L'$,
\item
for $X : \UU$, the canonical map $L'X \to LX$ is invertible.
\end{itemize}

When considering diagrams in a poset, we will restrict our attention to shapes that are themselves posets.
Note that in any poset, the limit of a (poset\hyp{}indexed) diagram coincides with the meet over the object components of the diagram.
Nonetheless, it is useful to speak about limits of diagrams because this allows us to constrain the relations between the inputs objects.

A poset $I$ is \emph{filtered} if it is merely inhabited and for any two elements $x_0, x_1 : I$, there merely exists $y : I$ with $x_0, x_1 \leq y$.
It is \emph{cofiltered} if $I^\op$ is filtered.
The following statement generalizes \cref{filter} to modalities.

\begin{proposition} \label{cofiltered-limit-subuniverses}
Let $(M, M') : I \to \Modality(\UU < \UU')$ be a $\UU$\hyp{}small cofiltered diagram.
If $\UU_{M_i}$ is $M'_i$\hyp{}modal for all $i : I$, then $\UU(M) : \UU'$ belongs to $\UU'(M')$.
\end{proposition}

\begin{proof}
Given $i : I$, we have to show that $\UU(M)$ is $M'_i$\hyp{}modal.
Because $I$ is cofiltered, we have
\[
\UU(M) = \UU\parens[\big]{(M_j)_{j \leq i}}
,\]
so it suffices to show that $\UU\parens[\big]{(M_j)_{j \leq i}}$ is $M'_i$\hyp{}modal.
By assumption, $\UU_{M_j}$ is $M'_j$\hyp{}modal, hence $M'_i$\hyp{}modal for $j \leq i$.
We now use that a type $X$ over $\UU_{M_i}$ ($M'_i$\hyp{}modal) is $M'_i$\hyp{}modal exactly if the map $X \to \UU_{M_i}$ is $M'_i$\hyp{}modal.
Given that $\UU_{M_j} \to \UU_{M_i}$ is $M'_i$\hyp{}modal for $j \leq i$, it suffices to show that $\UU\parens[\big]{(M_j)_{j \leq i}} \to \UU_{M_i}$ is $M'_i$\hyp{}modal.
Observe that the fibers of the latter embedding are products of the fibers of the former embeddings.
So the claim holds since modal types are closed under product (\cite[Lemma~1.26]{RijkeSS17}).
\end{proof}

\begin{corollary} \label{cofiltered-limit-universes-of-modals}
Let $(M, M') : I \to \Modality(\UU < \UU')$ be a $\UU$\hyp{}small cofiltered diagram with a structural meet $(\bigwedge M, \bigwedge M')$.
If $\UU_{M_i}$ is $M'_i$\hyp{}modal for all $i : I$, then $\UU_{\bigwedge M}$ is $\bigwedge M'$\hyp{}modal.
\end{corollary}

\begin{proof}
This is a direct consequence of \cref{cofiltered-limit-subuniverses} and the definition of structural meet.
\end{proof}

The following statement says that, up to the same size issues of \cref{modality-lex-via-universe}, lex modalities are closed under structural cofiltered limits of modalities.
In particular, structural cofiltered limits of ``size\hyp{}polymorphic'' modalities whose actions on maps preserve smallness of fibers preserve left exactness.

\begin{corollary} \label{cofiltered-limit-lex}
In the situation of \cref{cofiltered-limit-universes-of-modals}, if $M'_i$ is lex for $i : I$ and preserves maps with $\UU$\hyp{}small fibers, then $\bigwedge M$ is lex and $\bigwedge M'$ preserves maps with $\UU$\hyp{}small fibers.
\end{corollary}

\begin{proof}
This is the combination of \cref{modality-lex-via-universe} and \cref{cofiltered-limit-universes-of-modals}.
\end{proof}

Finally, we specialize to the important case of accessible modalities.

\begin{corollary} \label{cofiltered-limit-lex-accessible}
Let $M : I \to \Modality(\UU)$ be a $\UU$\hyp{}small cofiltered diagram.
If $M_i$ is lex and accessible for all $i : I$, then the meet $\bigwedge M$ exists and also has these properties.
\end{corollary}

\begin{proof}
Let $\UU'$ be a universe containing $\UU$.
Let $(M, M') : I \to \Modality(\UU < \UU')$ be the extension of $M$ given by~\cite[Theorem~3.36]{RijkeSS17}.
By~\cite[Theorem~3.29]{RijkeSS17}, the meet of $(M, M')$ exists, is structural, and $\bigwedge M$ is again accessible.
By~\cite[Theorem~3.11]{RijkeSS17}, $\UU_{M_i}$ is $M'_i$\hyp{}modal for $i : I$.
Applying \cref{cofiltered-limit-universes-of-modals}, $\UU_{\bigwedge M}$ is $\bigwedge M'$\hyp{}modal.
By~\cite[Theorem~3.29]{RijkeSS17}, this makes $\bigwedge M$ is lex.
\end{proof}

\begin{thebibliography}{10}

\bibitem{Aczel98}
Peter Aczel.
\newblock On relating type theories and set theories.
\newblock In Thorsten Altenkirch, Wolfgang Naraschewski, and Bernhard Reus,
  editors, {\em Types for Proofs and Programs, International Workshop {TYPES}
  '98, Kloster Irsee, Germany, March 27-31, 1998, Selected Papers}, volume 1657
  of {\em Lecture Notes in Computer Science}, pages 1--18. Springer, 1998.

\bibitem{AngiuliBCHHL}
Carlo Angiuli, Guillaume Brunerie, Thierry Coquand, Kuen-Bang {Hou (Favonia)},
  Robert Harper, and Daniel~R. Licata.
\newblock Cartesian cubical type theory.
\newblock Draft, December 2017.

\bibitem{AnnenkovCKS17}
Danil Annenkov, Paolo Capriotti, Nicolai Kraus, and Christian Sattler.
\newblock Two-level type theory and applications.
\newblock {\em CoRR}, abs/1705.03307, 2017.

\bibitem{AvigadKL15}
Jeremy Avigad, Krzysztof Kapulkin, and Peter~LeFanu Lumsdaine.
\newblock Homotopy limits in type theory.
\newblock {\em Math. Struct. Comput. Sci.}, 25(5):1040--1070, 2015.

\bibitem{Barr74}
Michael Barr.
\newblock Toposes without points.
\newblock {\em J. Pure Appl. Algebra}, 5:265--280, 1974.

\bibitem{BergnerR13}
Julia~E. Bergner and Charles Rezk.
\newblock Reedy categories and the {$\varTheta$}-construction.
\newblock {\em Math. Z.}, 274(1-2):499--514, 2013.

\bibitem{Beth56}
Evert~W. Beth.
\newblock {\em Semantic Construction of Intuitionistic Logic}.
\newblock Medededlingen der koninklijke Nederlandse Akademie van Wetenschappen,
  afd. Letterkunde. Nieuwe Reeks, Deel 19, No. 11. N. V. Noord-Hollandsche
  Uitgevers Maatschappij, Amsterdam, 1956.

\bibitem{Boulier18}
Simon~Pierre Boulier.
\newblock {\em Extending type theory with syntactic models. (Etendre la
  th\'{e}orie des types \`{a} l'aide de mod\`{e}les syntaxiques)}.
\newblock PhD thesis, Ecole nationale sup\'{e}rieure Mines-T\'{e}l\'{e}com
  Atlantique Bretagne Pays de la Loire, France, 2018.

\bibitem{CCHM15}
Cyril Cohen, Thierry Coquand, Simon Huber, and Anders M\"{o}rtberg.
\newblock Cubical type theory: {A} constructive interpretation of the
  univalence axiom.
\newblock In Tarmo Uustalu, editor, {\em 21st International Conference on Types
  for Proofs and Programs, {TYPES} 2015, May 18-21, 2015, Tallinn, Estonia},
  volume~69 of {\em LIPIcs}, pages 5:1--5:34. Schloss Dagstuhl -
  Leibniz-Zentrum fuer Informatik, 2015.

\bibitem{CHM18}
Thierry Coquand, Simon Huber, and Anders M\"{o}rtberg.
\newblock On higher inductive types in cubical type theory.
\newblock In Anuj Dawar and Erich Gr\"{a}del, editors, {\em Proceedings of the
  33rd Annual {ACM/IEEE} Symposium on Logic in Computer Science, {LICS} 2018,
  Oxford, UK, July 09-12, 2018}, pages 255--264. {ACM}, 2018.

\bibitem{CoquandMR17}
Thierry Coquand, Bassel Mannaa, and Fabian Ruch.
\newblock Stack semantics of type theory.
\newblock In {\em 32nd Annual {ACM/IEEE} Symposium on Logic in Computer
  Science, {LICS} 2017, Reykjavik, Iceland, June 20-23, 2017}, pages 1--11.
  {IEEE} Computer Society, 2017.

\bibitem{CoquandP88}
Thierry Coquand and Christine Paulin.
\newblock Inductively defined types.
\newblock In Per Martin{-}L\"{o}f and Grigori Mints, editors, {\em COLOG-88,
  International Conference on Computer Logic, Tallinn, USSR, December 1988,
  Proceedings}, volume 417 of {\em Lecture Notes in Computer Science}, pages
  50--66. Springer, 1988.

\bibitem{Dybjer95}
Peter Dybjer.
\newblock Internal type theory.
\newblock In Stefano Berardi and Mario Coppo, editors, {\em Types for Proofs
  and Programs, International Workshop TYPES'95, Torino, Italy, June 5-8, 1995,
  Selected Papers}, volume 1158 of {\em Lecture Notes in Computer Science},
  pages 120--134. Springer, 1995.

\bibitem{EilenbergZ50}
Samuel Eilenberg and J.~A. Zilber.
\newblock Semi-simplicial complexes and singular homology.
\newblock {\em Ann. of Math. (2)}, 51:499--513, 1950.

\bibitem{EGAI}
A.~Grothendieck.
\newblock \'{E}l\'{e}ments de g\'{e}om\'{e}trie alg\'{e}brique. {I}. {L}e
  langage des sch\'{e}mas.
\newblock {\em Inst. Hautes \'{E}tudes Sci. Publ. Math.}, 4:228, 1960.

\bibitem{Hofmann97}
Martin Hofmann.
\newblock Syntax and semantics of dependent types.
\newblock In {\em Semantics and logics of computation ({C}ambridge, 1995)},
  volume~14 of {\em Publ. Newton Inst.}, pages 79--130. Cambridge Univ. Press,
  Cambridge, 1997.

\bibitem{Joyal84}
Andr{\'{e}} Joyal.
\newblock Lettre {\`{a}} {Grothendieck}, 1984.

\bibitem{KaposiHS19}
Ambrus Kaposi, Simon Huber, and Christian Sattler.
\newblock Gluing for type theory.
\newblock In Herman Geuvers, editor, {\em 4th International Conference on
  Formal Structures for Computation and Deduction, {FSCD} 2019, June 24-30,
  2019, Dortmund, Germany}, volume 131 of {\em LIPIcs}, pages 25:1--25:19.
  Schloss Dagstuhl - Leibniz-Zentrum f\"{u}r Informatik, 2019.

\bibitem{Kraus15}
Nicolai Kraus.
\newblock {\em Truncation levels in homotopy type theory}.
\newblock PhD thesis, University of Nottingham, {UK}, 2015.

\bibitem{Kripke65}
Saul~A. Kripke.
\newblock Semantical analysis of intuitionistic logic. {I}.
\newblock In {\em Formal {S}ystems and {R}ecursive {F}unctions ({P}roc.
  {E}ighth {L}ogic {C}olloq., {O}xford, 1963)}, pages 92--130. North-Holland,
  Amsterdam, 1965.

\bibitem{MannaaC13}
Bassel Mannaa and Thierry Coquand.
\newblock Dynamic {Newton}-{Puiseux} theorem.
\newblock {\em J. Logic {\&} Analysis}, 5, 2013.

\bibitem{OrtonP16}
Ian Orton and Andrew~M. Pitts.
\newblock Axioms for modelling cubical type theory in a topos.
\newblock In Jean{-}Marc Talbot and Laurent Regnier, editors, {\em 25th {EACSL}
  Annual Conference on Computer Science Logic, {CSL} 2016, August 29 -
  September 1, 2016, Marseille, France}, volume~62 of {\em LIPIcs}, pages
  24:1--24:19. Schloss Dagstuhl - Leibniz-Zentrum fuer Informatik, 2016.

\bibitem{Quirin16}
Kevin Quirin.
\newblock {\em {Lawvere}-{Tierney} sheafification in Homotopy Type Theory.
  (Faisceautisation de {Lawvere}-{Tierney} en th\'{e}orie des types
  homotopiques)}.
\newblock PhD thesis, {\'{E}}cole des mines de Nantes, France, 2016.

\bibitem{RijkeSS17}
Egbert Rijke, Michael Shulman, and Bas Spitters.
\newblock Modalities in homotopy type theory.
\newblock {\em CoRR}, abs/1706.07526, 2017.

\bibitem{Sattler17}
Christian Sattler.
\newblock The equivalence extension property and model structures.
\newblock {\em CoRR}, abs/1704.06911, 2017.

\bibitem{SchreiberS14}
Urs Schreiber and Michael Shulman.
\newblock Quantum gauge field theory in cohesive homotopy type theory.
\newblock In Ross Duncan and Prakash Panangaden, editors, {\em Proceedings 9th
  Workshop on Quantum Physics and Logic, {QPL} 2012, Brussels, Belgium, 10-12
  October 2012}, volume 158 of {\em {EPTCS}}, pages 109--126, 2012.

\bibitem{Scott80}
Dana~S. Scott.
\newblock Relating theories of the {$\lambda $}-calculus.
\newblock In {\em To {H}. {B}. {C}urry: essays on combinatory logic, lambda
  calculus and formalism}, pages 403--450. Academic Press, London-New York,
  1980.

\bibitem{Shulman13}
Michael Shulman.
\newblock The univalence axiom for elegant reedy presheaves, 2013.

\bibitem{Shulman15}
Michael Shulman.
\newblock Univalence for inverse diagrams and homotopy canonicity.
\newblock {\em Mathematical Structures in Computer Science}, 25(5):1203--1277,
  2015.

\bibitem{Shulman18}
Michael Shulman.
\newblock {Brouwer}'s fixed-point theorem in real-cohesive homotopy type
  theory.
\newblock {\em Mathematical Structures in Computer Science}, 28(6):856--941,
  2018.

\bibitem{Shulman19}
Michael Shulman.
\newblock All $(\infty,1)$-toposes have strict univalent universes.
\newblock {\em CoRR}, abs/1904.07004, 2019.

\bibitem{SwanU19}
Andrew Swan and Taichi Uemura.
\newblock On {Church}'s thesis in cubical assemblies.
\newblock {\em CoRR}, abs/1905.03014, 2019.

\bibitem{TroelstraD88b}
A.~S. Troelstra and D.~van Dalen.
\newblock {\em Constructivism in mathematics. {V}ol. {II}}, volume 123 of {\em
  Studies in Logic and the Foundations of Mathematics}.
\newblock North-Holland Publishing Co., Amsterdam, 1988.
\newblock An introduction.

\bibitem{HoTT}
The {Univalent Foundations Program}.
\newblock {\em Homotopy Type Theory: Univalent Foundations of Mathematics}.
\newblock Institute for Advanced Study, 2013.

\bibitem{Voevodsky15}
Vladimir Voevodsky.
\newblock An experimental library of formalized mathematics based on the
  univalent foundations.
\newblock {\em Mathematical Structures in Computer Science}, 25(5):1278--1294,
  2015.

\bibitem{WeaverL20}
Matthew~Z. Weaver and Daniel~R. Licata.
\newblock A constructive model of directed univalence in bicubical sets.
\newblock In Holger Hermanns, Lijun Zhang, Naoki Kobayashi, and Dale Miller,
  editors, {\em {LICS} '20: 35th Annual {ACM/IEEE} Symposium on Logic in
  Computer Science, Saarbr\"{u}cken, Germany, July 8-11, 2020}, pages 915--928.
  {ACM}, 2020.

\bibitem{Wellen17}
Felix Wellen.
\newblock {\em Formalizing {C}artan Geometry in Modal Homotopy Type Theory}.
\newblock PhD thesis, Karlsruher Institut für Technologie, Germany, 2017.

\end{thebibliography}

%\input{lex3.tex}

\end{document}
